%% toc: P2 10 10.1,10.2,10.5,10.6
\section{Lindstr\"om's theorem for languages with function symbols, and proof routes for Barwise's survey claims}
\label{sec:p02-w02c-p02-003-general-lindstrom}

Barwise's languages contain function, relation and constant symbols, with
no restriction on the size of the set of symbols
\cite[ch.~A.1, \S3, p.~17]{Barwise1977FirstOrder}.  The form of his
Lindstr\"om theorem 5.7 proved in \cref{sec:p02-w02b-p02-001-lindstrom}
covers relational signatures only, and the chapter leaves both the
precise notion of a logic and the proof of 5.7 to his paper
\cite{Barwise1974Axioms}.  His chapter also reports a number of results
about particular theories (\S2, p.~16), about set-theoretic models (\S3,
p.~18) and about logics beyond first order (\S5, pp.~43--45) without
proof.

\emph{Purpose.}  This section carries 5.7 over to Barwise's general
languages, and for each of the reported results either writes a proof
from exact hypotheses or states precisely which steps of a proof are not
supplied in the handbook.  \emph{Inputs.}  The relational form
\cref{thm:p02-w02b-p02-001-barwise-lindstrom}; Barwise's notions of
structure, formula and satisfaction
(\cref{def:p02-w02-p02-005-bw-structure,def:p02-w02-p02-005-bw-syntax,def:p02-w02-p02-005-bw-satisfaction});
weak second-order logic, \(L_{\omega_1\omega}\), \(L(Q_0)\) and \(L(Q)\)
as in \cref{def:p02-w02-p02-005-bw-weak-so,def:p02-w02-p02-005-bw-infinitary};
the upward L\"owenheim--Skolem theorem \cref{thm:upward-ls}; the fact
that true arithmetic is not computably enumerable,
\cref{thm:p02-w02b-p02-034-ta-not-ce}; and G\"odel's second
incompleteness theorem \cref{thm:godel-second}.  \emph{Conclusions.}
\Cref{thm:p02-w02c-p02-003-general-lindstrom} (5.7 for general languages);
\cref{thm:p02-w02c-p02-003-b74-lindstrom,thm:p02-w02c-p02-003-b74-five-seven}
(Lindstr\"om's theorem and 5.7 for logics in the sense of Barwise's paper
\cite{Barwise1974Axioms}, to which the chapter refers);
\cref{thm:p02-w02c-p02-003-or-not,thm:p02-w02c-p02-003-small-model}
(Barwise's remarks at p.~44 on connectives and at p.~18 on small
set-theoretic models);
\cref{thm:p02-w02c-p02-003-infinitary-rule,thm:p02-w02c-p02-003-wso-translation,thm:p02-w02c-p02-003-hf,thm:p02-w02c-p02-003-beth}
(four of the claims of \S5; the first is completed by
\cref{thm:p02-w02c-p02-003-infinitary-completeness}, completeness of the
infinitary calculus \(\mathbf H_{\infty}\) of
\cref{def:p02-w02c-p02-003-hinf} for countable theories; the last through the interpolation
theorem \cref{thm:p02-w02c-p02-003-interpolation}, which is proved here
from a model existence theorem,
\cref{thm:p02-w02c-p02-003-model-existence});
\cref{thm:p02-w02c-p02-003-hf-recursion,thm:p02-w02c-p02-003-wso-sentences}
(Barwise's claims at p.~43 that functions can be defined by recursion in
\(\mathrm{HF}(\mathfrak M)\) and that his \S2 sentences are weak
second-order);
\cref{thm:p02-w02c-p02-003-wso-beth-failure} (Barwise's claim at p.~44
that Beth's theorem fails for weak second-order logic), through
\cref{thm:p02-w02c-p02-003-wso-n,thm:p02-w02c-p02-003-hf-arith,thm:p02-w02c-p02-003-truth-clauses};
\cref{thm:p02-w02c-p02-003-q0-equivalence} (Barwise's claim at p.~44 that
\(L(Q_0)\), \(\omega\)-logic and weak second-order logic are essentially
equivalent, in the sense of \cref{def:p02-w02c-p02-003-rpc});
\cref{thm:p02-w02c-p02-003-acf,thm:p02-w02c-p02-003-integers,thm:p02-w02c-p02-003-integers-second}
(four of the claims of p.~16, through Steinitz's classification
\cref{thm:p02-w02c-p02-003-steinitz}, the fundamental theorem of algebra
\cref{thm:p02-w02c-p02-003-fta}, Lagrange's four-square theorem
\cref{thm:p02-w02c-p02-003-lagrange} and the coding lemma
\cref{thm:p02-w02c-p02-003-ring-coding}; the last through the second
incompleteness theorem \cref{thm:godel-second}, applied with the provability
formula of \cref{def:p02-w02c-p02-003-sharp-formula}, whose derivability
conditions are \cref{thm:p02-w02c-p02-003-sharp-provability}, and the
explicit enumeration \cref{thm:p02-w02c-p02-003-sharp-enumeration});
\cref{thm:p02-w02c-p02-003-hf-sigma-recursion} (Makkai's
\(\Sigma\)-recursion on membership for \(\mathrm{HF}(\mathfrak M)\));
\cref{thm:p02-w02c-p02-003-keisler-completeness,thm:p02-w02c-p02-003-keisler-compactness}
(Keisler's Completeness and Compactness Theorems for \(L(Q)\), reported by
Barwise at pp.~44--45); and the two boundary statements
\cref{def:p02-w02c-p02-003-beyond-boundary,def:p02-w02c-p02-003-theories-boundary}
for the claims whose proofs are not written here.

\subsection{Function and constant symbols}
\label{sec:p02-w02c-p02-003-functions}

The passage from relational signatures to Barwise's languages replaces
each function and constant symbol by a relation symbol for its graph.

\begin{hbdefinition}{def:p02-w02c-p02-003-relational-image}{Relational image of a language and of a structure}
Let \(L\) be a language in Barwise's sense: a set of relation, function and
constant symbols, each relation and function symbol with a positive arity
\cite[ch.~A.1, \S3, p.~17]{Barwise1977FirstOrder}.  Its \emph{relational
image} \(L^{r}\) contains every relation symbol of \(L\), an
\((n+1)\)-ary relation symbol \(R_f\) for each \(n\)-ary function symbol
\(f\in L\), and a unary relation symbol \(P_c\) for each constant symbol
\(c\in L\); the new symbols are distinct from each other and from the
symbols of \(L\).  For a structure \(\mathfrak A\) for \(L\), the structure
\(\mathfrak A^{r}\) for \(L^{r}\) has the domain of \(\mathfrak A\), the
same interpretation of each relation symbol of \(L\), and
\[
 R_f^{\mathfrak A^{r}}=\{(a_1,\ldots,a_n,b):f^{\mathfrak A}(a_1,\ldots,a_n)=b\},
 \qquad
 P_c^{\mathfrak A^{r}}=\{c^{\mathfrak A}\}.
\]
For \(L_0\subseteq L\), \(L_0^{r}\subseteq L^{r}\) and
\((\mathfrak A\restriction L_0)^{r}=\mathfrak A^{r}\restriction L_0^{r}\).
For a formula \(\chi\) of \(L^{r}\) in first-order logic or
\(L_{\omega_1\omega}\), \(\chi^{-}\) is the formula of \(L\) in the same
logic obtained by replacing each atomic subformula
\(R_f(x_1,\ldots,x_n,y)\) by \(f(x_1,\ldots,x_n)=y\) and each \(P_c(x)\) by
\(c=x\), and leaving all other symbols unchanged.  Since \(L^{r}\) has no
function or constant symbols, its terms are variables, so every atomic
formula of \(L^{r}\) has one of these forms, the form \(R(x_1,\ldots,x_n)\)
with \(R\in L\), or the form \(x=y\).
\end{hbdefinition}

For a term \(t\) of \(L\) and a variable \(y\) not occurring in \(t\),
define a first-order formula \(E_t(y)\) of \(L^{r}\) by recursion on
\(t\): \(E_x(y)\) is \(x=y\); \(E_c(y)\) is \(P_c(y)\); and
\(E_{f(t_1,\ldots,t_n)}(y)\) is
\(\exists z_1\cdots\exists z_n\,(E_{t_1}(z_1)\land\cdots\land
E_{t_n}(z_n)\land R_f(z_1,\ldots,z_n,y))\) with \(z_1,\ldots,z_n\) new.  For
a formula \(\chi\) of \(L\) in first-order logic, \(L_{\omega_1\omega}\),
\(L(Q_0)\) or \(L(Q)\), the \emph{unnesting} \(\chi^{u}\) is obtained by
replacing each atomic subformula \(R(t_1,\ldots,t_n)\) by
\(\exists z_1\cdots\exists z_n\,(E_{t_1}(z_1)\land\cdots\land
E_{t_n}(z_n)\land R(z_1,\ldots,z_n))\) and each atomic subformula \(t=t'\)
by \(\exists z\,(E_t(z)\land E_{t'}(z))\), with new variables, and leaving
connectives and quantifiers (including \(Q\), \(\bigwedge\), \(\bigvee\))
unchanged.

\begin{hblemma}{thm:p02-w02c-p02-003-translations}{proved=annex:p02-w02c-p02-003-translations}{Translations between a language and its relational image}
Let \(\mathfrak A\) be a structure for \(L\) and \(s\) an assignment in its
domain.
\begin{enumerate}
\item For every formula \(\chi\) of \(L^{r}\) in first-order logic or
      \(L_{\omega_1\omega}\), \(\mathfrak A\models\chi^{-}[s]\) iff
      \(\mathfrak A^{r}\models\chi[s]\).
\item For every term \(t\) of \(L\) and variable \(y\) not in \(t\),
      \(\mathfrak A^{r}\models E_t(y)[s]\) iff
      \(s(y)=t^{\mathfrak A}(s)\).
\item For every formula \(\chi\) of \(L\) in first-order logic,
      \(L_{\omega_1\omega}\), \(L(Q_0)\) or \(L(Q)\),
      \(\mathfrak A\models\chi[s]\) iff \(\mathfrak A^{r}\models\chi^{u}[s]\),
      and \(\chi^{u}\) is a formula of \(L^{r}\) in the same logic.
\end{enumerate}
\end{hblemma}

\begin{hbconvention}{def:p02-w02c-p02-003-general-reading}{Reading 5.7 for Barwise's languages}
A \emph{logic on Barwise's languages} \(\mathcal L^{\ast}\) assigns to each
language \(L\) (function, relation and constant symbols, of any
cardinality) a class \(\mathcal L^{\ast}(L)\) of sentences and a relation
\(\mathfrak A\models\varphi\) between structures for \(L\) and members of
\(\mathcal L^{\ast}(L)\).  It is \emph{regular} if:
\begin{enumerate}
\item its restriction to relational languages is a regular abstract logic
      in the relational sense of \cref{sec:p02-fol-metatheorem-summary};
\item (\emph{occurrence}) if \(L_0\subseteq L\) then
      \(\mathcal L^{\ast}(L_0)\subseteq\mathcal L^{\ast}(L)\) and
      \(\mathfrak A\models\varphi\) iff
      \(\mathfrak A\restriction L_0\models\varphi\) for
      \(\varphi\in\mathcal L^{\ast}(L_0)\); and each
      \(\varphi\in\mathcal L^{\ast}(L)\) lies in
      \(\mathcal L^{\ast}(L_0)\) for some finite \(L_0\subseteq L\), its
      \emph{occurrence language};
\item (\emph{replacement}) for every \(L\) and
      \(\varphi\in\mathcal L^{\ast}(L)\) there is
      \(\varphi^{r}\in\mathcal L^{\ast}(L^{r})\) such that
      \(\mathfrak A\models\varphi\) iff \(\mathfrak A^{r}\models\varphi^{r}\)
      for every structure \(\mathfrak A\) for \(L\);
\item every first-order sentence of \(L\) belongs to
      \(\mathcal L^{\ast}(L)\), with its usual satisfaction relation.
\end{enumerate}
\(\mathcal L^{\ast}\) \emph{satisfies the Compactness Theorem} and
\emph{the L\"owenheim--Skolem theorem} if parts~1 and~2 of
\cref{def:p02-w02b-p02-001-lindstrom-reading} hold with ``every relational
signature'' replaced by ``every language''.  \emph{First-order logic is
the only such logic} means that every \(\varphi\in\mathcal L^{\ast}(L)\)
with occurrence language \(L_0\) has a first-order sentence \(\psi\) of
\(L_0\), possibly containing function and constant symbols, with
\(\mathfrak A\models\varphi\) iff \(\mathfrak A\models\psi\) for every
structure \(\mathfrak A\) for \(L\).
\end{hbconvention}

By part~3 of \cref{thm:p02-w02c-p02-003-translations}, first-order logic,
\(L_{\omega_1\omega}\), \(L(Q_0)\) and \(L(Q)\) meet the replacement
condition, with \(\varphi^{r}=\varphi^{u}\).

\begin{hbtheorem}{thm:p02-w02c-p02-003-general-lindstrom}{proved=annex:p02-w02c-p02-003-general-lindstrom}{Barwise's 5.7 for languages with function and constant symbols}
Let \(\mathcal L^{\ast}\) be a regular logic on Barwise's languages in the
sense of \cref{def:p02-w02c-p02-003-general-reading} that satisfies the
Compactness Theorem and the L\"owenheim--Skolem theorem in that sense.
Then for every language \(L\) and every \(\varphi\in\mathcal L^{\ast}(L)\)
with occurrence language \(L_0\) there is a first-order sentence \(\psi\)
of \(L_0\) such that \(\mathfrak A\models\varphi\) iff
\(\mathfrak A\models\psi\) for every structure \(\mathfrak A\) for \(L\).
Conversely every first-order sentence of \(L\) belongs to
\(\mathcal L^{\ast}(L)\).  The proof uses the Compactness and
L\"owenheim--Skolem hypotheses only for relational languages.
\end{hbtheorem}

\begin{hbexample}{def:p02-w02c-p02-003-replacement-example}{Replacement for a sentence of \(L(Q)\) with a function symbol}
Let \(L=\{f,c\}\) with \(f\) unary and \(c\) a constant symbol, and let
\(\varphi\) be the \(L(Q)\)-sentence \(Qx\,(f(x)=c)\): uncountably many
elements are mapped by \(f\) to \(c\).  Then \(E_{f(x)}(z)\) is
\(\exists z_1\,(z_1=x\land R_f(z_1,z))\) and \(E_c(z)\) is \(P_c(z)\), so
\[
 \varphi^{u}=Qx\,\exists z\,\bigl(\exists z_1\,(z_1=x\land R_f(z_1,z))\land P_c(z)\bigr).
\]
In \(\mathfrak A^{r}\), for a fixed \(x\) the inner formula holds iff the
value \(z\) with \(R_f(x,z)\), namely \(f^{\mathfrak A}(x)\), is the unique
member \(c^{\mathfrak A}\) of \(P_c\); so \(\mathfrak A^{r}\models\varphi^{u}\)
iff uncountably many \(x\) have \(f^{\mathfrak A}(x)=c^{\mathfrak A}\), as
part~3 of \cref{thm:p02-w02c-p02-003-translations} requires.  The theorem
does not apply to \(L(Q)\), which fails the L\"owenheim--Skolem theorem
(part~4 of \cref{thm:p02-w02-p02-005-bw-extended-ls}); in the relational
image the failure is witnessed by \(Qx\,(x=x)\), which also has no
first-order equivalent (\cref{thm:p02-w02b-p02-001-lindstrom-sharp}).
\end{hbexample}

\begin{hbwarning}{def:p02-w02c-p02-003-general-limit}{What the general reading adds}
The occurrence and replacement conditions of
\cref{def:p02-w02c-p02-003-general-reading} are not named in 5.7, which
mentions only closure under \(\land\), \(\neg\), \(\exists\); the chapter
leaves the precise notion of a logic to \cite{Barwise1974Axioms}.  That
notion is set out in \cref{sec:p02-w02c-p02-003-b74}, the two readings are
compared item by item in \cref{def:p02-w02c-p02-003-b74-comparison}, and
5.7 is proved under Barwise's own axioms, without any replacement
condition, as \cref{thm:p02-w02c-p02-003-b74-lindstrom}.  For a logic that
meets the conditions of \cref{def:p02-w02c-p02-003-general-reading} other
than replacement but not Barwise's translation axiom, the replacement
condition is what ties the sentences of a language with function symbols to
those of a relational language, and whether
\cref{thm:p02-w02c-p02-003-general-lindstrom} holds for such a logic is not
settled here.  The theorem inherits every qualification of the relational
form recorded in \cref{def:p02-w02b-p02-001-lindstrom-limit}.
\end{hbwarning}

\subsection{Barwise's axioms for a logic, and 5.7 under them}
\label{sec:p02-w02c-p02-003-b74}

The chapter sends the reader to Barwise's paper ``Axioms for abstract
model theory'' both for the precise notion of a logic and for a proof of
5.7 \cite[ch.~A.1, \S5, 5.6--5.7, p.~45]{Barwise1977FirstOrder}.  The
paper states its axioms in Part~I, \S7, and proves Lindstr\"om's theorem in
Part~II, \S3, as Corollary~3.4 \cite[Part~II, \S3, 3.4,
p.~259]{Barwise1974Axioms}.

\emph{Purpose.}  This subsection states the paper's notion of a logic
exactly, records how it differs from the readings used in
\cref{def:p02-w02b-p02-001-lindstrom-reading} and
\cref{def:p02-w02c-p02-003-general-reading}, and proves the paper's
Corollary~3.4, and with it 5.7, for every logic in the paper's sense.
\emph{Inputs.}  The paper's definitions of languages, interpretations,
pulled-back structures, categories of languages and logics; first-order
satisfaction and the relational image
(\cref{def:p02-w02c-p02-003-relational-image}) with part~1 of
\cref{thm:p02-w02c-p02-003-translations}; the coincidence lemma
\cref{lem:fol-coincidence}.  \emph{Conclusion.}
\Cref{thm:p02-w02c-p02-003-b74-lindstrom}: a logic in the paper's sense
with the L\"owenheim property and either the upward L\"owenheim--Skolem
property or countable compactness is equivalent to first-order logic, for
languages with relation, function and constant symbols of any size.

\begin{hbdefinition}{def:p02-w02c-p02-003-b74-languages}{Languages, interpretations and pulled-back structures (Barwise 1974, Part~I, \S\S3--6)}
A fixed unary relation symbol \(\mathbf V\) is used throughout; it names the
range of the quantifiers.
\begin{enumerate}
\item A \emph{language} is a set \(L\) of relation, function and constant
      symbols with \(\mathbf V\in L\); each relation and function symbol has
      a finite number of argument places, and a constant symbol is a function
      symbol with \(0\) arguments.  A \emph{partial structure}
      \(\mathfrak M\) for \(L\) assigns to \(\mathbf V\) a set \(M\), its
      universe, to each \(n\)-ary relation symbol \(R\) a set
      \(R^{\mathfrak M}\subseteq M^{n}\), and to each \(n\)-ary function
      symbol \(f\) a partial function from \(M^{n}\) to \(M\); for a constant
      symbol \(c\), \(c^{\mathfrak M}\) is an element of \(M\) if it is
      defined.  It is a \emph{structure for \(L\)}, or \(L\)-structure, when
      every \(f^{\mathfrak M}\) is total and every \(c^{\mathfrak M}\) is
      defined \cite[Part~I, \S3, p.~228]{Barwise1974Axioms}.  The paper does
      not require \(M\) to be nonempty.
\item Let \(k\geq0\).  A \emph{\(k\)-term interpretation}
      \(\alpha:L\to K\) consists of terms \(t_1^{\alpha},\ldots,t_k^{\alpha}\)
      of \(K\) and a map \(\alpha\) from \(L\) to \(K\) that sends each
      \(n\)-ary relation symbol to an \((n+k)\)-ary relation symbol and each
      \(n\)-ary function symbol to an \((n+k)\)-ary function symbol, and
      sends no symbol other than \(\mathbf V\) to \(\mathbf V\).  We write
      \(s^{\alpha}\) for \(\alpha(s)\).  A \(0\)-term interpretation with
      \(\mathbf V^{\alpha}=\mathbf V\) is \emph{simple}; a one-one simple
      interpretation is an \emph{embedding} \cite[Part~I, \S4, 4.1--4.2,
      pp.~229--231]{Barwise1974Axioms}.  The category \(\mathcal C_\infty\)
      has all languages as objects and all interpretations as morphisms.
      For an inclusion \(i:L_0\subseteq L\) and an interpretation
      \(\alpha:L\to K\), the composite \(\alpha i\) is the restriction of
      \(\alpha\) to \(L_0\), with the same terms.
\item For \(\alpha:L\to K\) with terms \(t_1,\ldots,t_k\) and a
      \(K\)-structure \(\mathfrak N\), let \(a_j=t_j^{\mathfrak N}\).  The
      \emph{pull-back} \(\mathfrak N^{-\alpha}\) is the partial structure for
      \(L\) with universe
      \(N_\alpha=\{b:(b,a_1,\ldots,a_k)\in(\mathbf V^{\alpha})^{\mathfrak N}\}\),
      with \(R^{\mathfrak N^{-\alpha}}=\{\bar b\in N_\alpha^{n}:(\bar b,\bar a)\in(R^{\alpha})^{\mathfrak N}\}\),
      and with \(f^{\mathfrak N^{-\alpha}}\) the restriction to
      \(N_\alpha^{n}\) of \(\bar b\mapsto(f^{\alpha})^{\mathfrak N}(\bar b,\bar a)\),
      a value outside \(N_\alpha\) counting as undefined.  \(\mathfrak N\) is
      \emph{\(\alpha\)-invertible} when \(\mathfrak N^{-\alpha}\) is an
      \(L\)-structure, that is, when \(N_\alpha\) is closed under these
      functions and contains the value of each constant symbol
      \cite[Part~I, \S5, p.~231]{Barwise1974Axioms}.  For an inclusion the
      pull-back is the reduct; for \(\alpha:L\to L\cup\{U\}\) that sends
      \(\mathbf V\) to a new unary \(U\) and fixes the other symbols, it is
      the relativization to \(U\) when \(U\) is closed under the functions
      of \(L\) \cite[Part~I, \S5, 5.1--5.2, p.~231]{Barwise1974Axioms}.
\item A \emph{category of languages} is a subcategory \(\mathcal C\) of
      \(\mathcal C_\infty\) such that: (i)~\(L\cap K\) and \(L\cup K\) are
      in \(\mathcal C\) when \(L,K\) are; (ii)~for \(L\subseteq K\) in
      \(\mathcal C\) the inclusion is a morphism of \(\mathcal C\);
      (iii)~every finite language is isomorphic to one in \(\mathcal C\),
      and, for each \(K\) in \(\mathcal C\), to one in \(\mathcal C\) whose
      only common symbol with \(K\) is \(\mathbf V\); (iv)~a \(0\)-term
      interpretation between languages of \(\mathcal C\) that is the
      identity except on finitely many symbols is a morphism of
      \(\mathcal C\); (v)~for each \(L\) in \(\mathcal C\) there are a
      constant symbol \(d\notin L\), a language \(K\) in \(\mathcal C\) with
      \(d\in K\), and a one-one \(1\)-term interpretation \(L\to K\) with term
      \(d\) that is a morphism of \(\mathcal C\).  The paper calls this
      definition tentative \cite[Part~I, \S6, 6.2, pp.~233--234]{Barwise1974Axioms}.
\end{enumerate}
\end{hbdefinition}

An isomorphism of \(\mathcal C_\infty\) is a one-one interpretation onto its
target \cite[Part~I, \S4, 4.3, p.~230]{Barwise1974Axioms}; since
\(\mathbf V\) is then a value, and only \(\mathbf V\) can be sent to
\(\mathbf V\), it is a \(0\)-term interpretation that fixes \(\mathbf V\)
and preserves the kind and arity of every symbol.  So clause~(iii)
supplies, inside \(\mathcal C\), any finite stock of new symbols of
prescribed kinds and arities, and clause~(i) adjoins them to a given
language.

\begin{hbdefinition}{def:p02-w02c-p02-003-b74-logic}{A logic in Barwise's sense, and the hypotheses of his Corollary~3.4}
A \emph{logic} \(\mathcal L^{\ast}\) on a category of languages
\(\mathcal C\) consists of a functor \(L\mapsto L^{\ast}\),
\(\alpha\mapsto\alpha^{\ast}\) from \(\mathcal C\) to classes, whose values
are the \emph{\(\mathcal L^{\ast}\)-sentences} of each language, and a
relation \(\mathfrak M\models\varphi\), holding only when \(\mathfrak M\) is
an \(L\)-structure for some \(L\) in \(\mathcal C\) with \(\varphi\in L^{\ast}\),
such that \cite[Part~I, \S7, pp.~234--235]{Barwise1974Axioms}:
\begin{enumerate}
\item (\emph{Occurrence Axiom}) each \(\mathcal L^{\ast}\)-sentence
      \(\varphi\) has a least language \(L_\varphi\) in \(\mathcal C\) with
      \(\varphi\in L_\varphi^{\ast}\), and for every inclusion morphism
      \(i:L_\varphi\subseteq K\) of \(\mathcal C\) the map \(i^{\ast}\) is the
      inclusion \(L_\varphi^{\ast}\subseteq K^{\ast}\);
\item (\emph{Isomorphism Axiom}) if \(\mathfrak M\models\varphi\) and
      \(\mathfrak M\cong\mathfrak N\), then \(\mathfrak N\models\varphi\);
\item (\emph{Translation Axiom}) for every \(\mathcal L^{\ast}\)-sentence
      \(\varphi\), every morphism \(\alpha:L_\varphi\to K\) of \(\mathcal C\)
      and every \(K\)-structure \(\mathfrak N\),
      \(\mathfrak N\models\varphi^{\alpha}\) iff \(\mathfrak N\) is
      \(\alpha\)-invertible and \(\mathfrak N^{-\alpha}\models\varphi\),
      where \(\varphi^{\alpha}=\alpha^{\ast}(\varphi)\).
\end{enumerate}
The hypotheses of Corollary~3.4 are the standing assumptions of Part~II,
\S\S2--3, and two properties:
\begin{enumerate}\setcounter{enumi}{3}
\item (\emph{atomic sentences}) for every \(L\) in \(\mathcal C\), every
      atomic sentence \(t_1=t_2\) or \(R(t_1,\ldots,t_n)\) of \(L\) (closed
      terms only) is in \(L^{\ast}\), with its usual truth condition
      \cite[Part~I, \S8, p.~237]{Barwise1974Axioms};
\item (\emph{conjunction and negation}) for every \(L\) there are
      operations \(L^{\ast}\times L^{\ast}\to L^{\ast}\) and
      \(L^{\ast}\to L^{\ast}\), commuting with the maps \(\alpha^{\ast}\) of
      embeddings, written \((\varphi\land\psi)\) and \(\neg\varphi\), with
      \(\mathfrak M\models(\varphi\land\psi)\) iff both hold and
      \(\mathfrak M\models\neg\varphi\) iff \(\mathfrak M\not\models\varphi\),
      for \(L\)-structures \(\mathfrak M\) \cite[Part~I, \S9, 9.1, 9.9,
      pp.~238, 242]{Barwise1974Axioms};
\item (\emph{existential quantifier}) for every \(L\), every constant
      symbol \(c\in L\) and every \(\varphi\in L^{\ast}\) there is a sentence
      \(q(c,\varphi)\in L^{\ast}\) in which \(c\) does not occur, the
      operation commuting with embeddings, such that
      \(\mathfrak M\models q(c,\varphi)\) iff some \(a\) in the universe has
      \((\mathfrak M,a)\models\varphi\), where \((\mathfrak M,a)\) is
      \(\mathfrak M\) with \(c\), added to its language if it is not there,
      interpreted by \(a\) \cite[Part~I, \S9,
      9.3, 9.7, pp.~239, 241]{Barwise1974Axioms};
\item (\emph{finite occurrence}) \(L_\varphi\) is finite for every
      \(\mathcal L^{\ast}\)-sentence \(\varphi\)
      \cite[Part~II, \S2, p.~254]{Barwise1974Axioms};
\item (\emph{L\"owenheim property}) every \(\mathcal L^{\ast}\)-sentence that
      has a model has one of cardinality at most \(\aleph_0\)
      \cite[Part~II, \S2, p.~255]{Barwise1974Axioms};
\end{enumerate}
together with one of
\begin{enumerate}\setcounter{enumi}{8}
\item[(a)] (\emph{upward L\"owenheim--Skolem}) every
      \(\mathcal L^{\ast}\)-sentence with an infinite model has an
      uncountable model;
\item[(b)] (\emph{countable compactness}) if \(T\) is a countable subset of
      \(L^{\ast}\) for some \(L\) in \(\mathcal C\) and every finite subset
      of \(T\) has a model, then \(T\) has a model
\end{enumerate}
\cite[Part~II, \S3, 3.4, p.~259]{Barwise1974Axioms}.  For logics
\(\mathcal L^{\ast}\), \(\mathcal L^{\#}\) on \(\mathcal C\),
\(\mathcal L^{\#}\) is \emph{as strong as} \(\mathcal L^{\ast}\) if for
every \(\mathcal L^{\ast}\)-sentence \(\varphi\) there is an
\(\mathcal L^{\#}\)-sentence \(\psi\) whose symbols all occur in \(\varphi\)
and which has the same \(L_\varphi\)-structures as models; the logics are
\emph{equivalent}, \(\mathcal L^{\ast}\equiv\mathcal L^{\#}\), if each is as
strong as the other \cite[Part~II, \S3, 3.1, p.~257]{Barwise1974Axioms}.
\(L_{\omega\omega}\) is first-order logic on the same category: its
sentences of \(L\) are the first-order sentences built from the symbols of
\(L\) with \(=\) as a logical symbol, and its quantifiers range over the
universe \cite[Part~I, \S1, p.~222]{Barwise1974Axioms}.
\end{hbdefinition}

\begin{hbtheorem}{thm:p02-w02c-p02-003-b74-lindstrom}{proved=annex:p02-w02c-p02-003-b74-lindstrom}{Lindstr\"om's theorem for logics in Barwise's sense (Barwise 1974, Corollary~3.4)}
Let \(\mathcal L^{\ast}\) be a logic on a category of languages
\(\mathcal C\) satisfying conditions 4--8 of
\cref{def:p02-w02c-p02-003-b74-logic}, and suppose it satisfies (a) or
(b).  Then \(\mathcal L^{\ast}\equiv L_{\omega\omega}\).  In detail:
\begin{enumerate}
\item for every \(\mathcal L^{\ast}\)-sentence \(\varphi\) there is a
      first-order sentence \(\psi\) of \(L_\varphi\) such that, for every
      \(L_\varphi\)-structure \(\mathfrak M\) (the empty one included, when
      it exists), \(\mathfrak M\models\varphi\) iff \(\mathfrak M\models\psi\);
\item for every \(L\) in \(\mathcal C\) and every first-order sentence
      \(\psi\) of \(L\) there is \(\theta\in L^{\ast}\) with the same
      \(L\)-structures as models; when \(L\) is the occurrence language of
      \(\psi\) in \(L_{\omega\omega}\), every symbol occurring in \(\theta\)
      occurs in \(\psi\).
\end{enumerate}
\end{hbtheorem}

\begin{hbcorollary}{thm:p02-w02c-p02-003-b74-five-seven}{proved=annex:p02-w02c-p02-003-b74-five-seven}{Barwise's 5.7 for logics in his sense}
Let \(\mathcal L^{\ast}\) be a logic on \(\mathcal C\) satisfying conditions
4--6 of \cref{def:p02-w02c-p02-003-b74-logic} (closure under \(\land\),
\(\neg\), \(\exists\), and the atomic sentences) and condition~7.  Suppose
\(\mathcal L^{\ast}\) satisfies the Compactness Theorem 2.4, read for every
\(L\) in \(\mathcal C\) and every set of \(\mathcal L^{\ast}\)-sentences of
\(L\), and the L\"owenheim--Skolem theorem 2.5 with \(\kappa=\aleph_0\) as
in part~2 of \cref{def:p02-w02b-p02-001-lindstrom-reading}.  Then
\(\mathcal L^{\ast}\equiv L_{\omega\omega}\).
\end{hbcorollary}

The proof of \cref{thm:p02-w02c-p02-003-b74-lindstrom} follows the plan of
\cref{thm:p03-lindstrom}: a back-and-forth system between two structures
is written into one sentence of the logic, the L\"owenheim property turns a
model of that sentence into an isomorphism, and either hypothesis~(b) or
hypothesis~(a) excludes a sentence that is not equivalent to a sentence of
bounded quantifier rank.  Three points differ.  Function and constant
symbols are handled through partial isomorphisms of relational images,
which preserve the graphs of the functions.  The two copies of a structure
are placed inside one structure by Barwise's translation axiom, applied to
a relativizing interpretation that also renames the symbols of the second
copy; the closure of each copy under its functions is stated as a
first-order axiom, since the translation of a sentence is false in a
structure that is not invertible.  Under hypothesis~(a), the sentences
built for different quantifier ranks are gathered into one sentence by a
\(1\)-term interpretation, Barwise's indexed union, and that sentence has
infinite models but no uncountable one.  Barwise's own route is different:
it passes through Karp's theorem, his Theorem~2.2 (the L\"owenheim property
gives the Karp property), Theorem~3.2 and Theorem~1.2
\cite[Part~II, \S\S1--3, pp.~249--259]{Barwise1974Axioms}; the proof here
does not use Karp's theorem, whose proof the paper refers to an earlier
survey \cite[Part~I, \S10, p.~243]{Barwise1974Axioms}.

\begin{hbwarning}{def:p02-w02c-p02-003-b74-comparison}{How Barwise's notion differs from the readings used for 5.7}
\begin{enumerate}
\item \emph{Morphisms.}  Barwise's translation axiom covers every
      interpretation of the category, including relativizations,
      renamings, the identification of two symbols and \(k\)-term
      interpretations such as indexed unions
      \cite[Part~I, \S5, 5.2--5.4, pp.~231--232]{Barwise1974Axioms}.  The
      relational regular logics of \cref{sec:p02-fol-metatheorem-summary}
      are required only to support renamings and relativization to a
      nonempty unary predicate, and only for relational signatures.
\item \emph{Function symbols.}  An interpretation sends function symbols
      to function symbols, so no morphism of Barwise's category replaces a
      function symbol by a relation symbol for its graph.  The replacement
      condition of \cref{def:p02-w02c-p02-003-general-reading} is not among
      Barwise's axioms, and \cref{thm:p02-w02c-p02-003-b74-lindstrom} does
      not use it; in its place the proof uses the translation axiom for
      languages that contain function symbols.
\item \emph{Structures.}  Barwise's structures need not have nonempty
      universes, and pull-backs can be partial structures; the chapter's
      structures have nonempty universes
      \cite[ch.~A.1, \S3, 3.1, p.~17]{Barwise1977FirstOrder}.
      \Cref{thm:p02-w02c-p02-003-b74-lindstrom} covers the empty structure,
      so its conclusion does not depend on how this point is settled.
\item \emph{Hypotheses.}  The chapter's 5.7 names the Compactness Theorem
      2.4, for arbitrary sets of sentences, and the L\"owenheim--Skolem
      theorem 2.5 with \(\kappa=\aleph_0\).  Barwise's Corollary~3.4 needs
      less: compactness for countable sets in one language, or instead the
      upward L\"owenheim--Skolem property, together with the L\"owenheim
      property for single sentences.  It also assumes finite occurrence
      languages, a standing assumption of Part~II, \S\S2--3 that Barwise
      lists as one he expects a better axiom to remove (problem~P5)
      \cite[Part~II, \S5, p.~264]{Barwise1974Axioms}; 5.7 does not mention
      it, and \cref{thm:p02-w02c-p02-003-b74-five-seven} keeps it as a
      hypothesis.
\item \emph{Languages.}  A logic in Barwise's sense lives on a category of
      languages, which may omit languages; the chapter's languages are
      unrestricted in size \cite[ch.~A.1, \S3, p.~17]{Barwise1977FirstOrder},
      which corresponds to \(\mathcal C=\mathcal C_\infty\).
\item \emph{Conclusion.}  Barwise's conclusion is the equivalence
      \(\mathcal L^{\ast}\equiv L_{\omega\omega}\) of \(3.1\): a first-order
      sentence in the symbols that occur, and the converse inclusion.
      \Cref{thm:p02-w02b-p02-001-barwise-lindstrom} and
      \cref{thm:p02-w02c-p02-003-general-lindstrom} give the first half
      over the occurrence signature and take the converse as part of their
      definitions.
\end{enumerate}
Barwise remarks that the 1974 approach differs from Lindstr\"om's own in
that its logics allow relativization and in that it treats sentences, not
only classes of structures, as the primary objects; by working with his
notion he removes exceptions that Lindstr\"om's formulation had to make for
finite models \cite[Part~II, \S3, p.~259]{Barwise1974Axioms}.
\end{hbwarning}

\begin{hbexample}{def:p02-w02c-p02-003-b74-example}{Relativizing a sentence with a function symbol}
Let \(L=\{\mathbf V,f\}\) with \(f\) unary, let \(\varphi\) be the
first-order sentence \(\forall x\,\neg f(x)=x\), and let
\(\alpha:L\to L\cup\{U\}\) send \(\mathbf V\) to a new unary \(U\) and fix
\(f\).  In the structure \(\mathfrak N\) with universe
\(\{0,1,2\}\), \(f(0)=1\), \(f(1)=0\), \(f(2)=2\) and \(U=\{0,1\}\), the set
\(U\) is closed under \(f\), the pull-back is the two-element structure
in which \(f\) swaps \(0\) and \(1\), and \(\mathfrak N\models\varphi^{\alpha}\).
With \(U=\{0\}\) the set is not closed under \(f\), so \(\mathfrak N\) is not
\(\alpha\)-invertible: \(\mathfrak N\not\models\varphi^{\alpha}\) and also
\(\mathfrak N\not\models(\neg\varphi)^{\alpha}\), while
\(\mathfrak N\models\neg(\varphi^{\alpha})\).  This is why the proof of
\cref{thm:p02-w02c-p02-003-b74-lindstrom} adds the closure of each copy as a
separate axiom, and why Barwise observes that \((\neg\varphi)^{\alpha}\) and
\(\neg(\varphi^{\alpha})\) agree only on invertible structures
\cite[Part~I, \S9, after 9.11, pp.~243--244]{Barwise1974Axioms}.
\end{hbexample}

\subsection{Connectives from \(\lor\) and \(\neg\), and small set-theoretic models}
\label{sec:p02-w02c-p02-003-or-not}

Barwise remarks that every finitary propositional connective is definable
from those of first-order logic, and in fact from \(\lor\) and \(\neg\)
\cite[ch.~A.1, \S5, 5.5, p.~44]{Barwise1977FirstOrder}; his paper states the
same for \(\neg\) and \(\land\), with no proof beyond calling the fact well known
\cite[Part~I, \S9, 9.10, p.~242]{Barwise1974Axioms}.  An \(n\)-ary
connective \(\#\) with truth function \(p:\{0,1\}^{n}\to\{0,1\}\) is added
to first-order logic by the clause: \(\mathfrak M\models\#(\varphi_1,\ldots,
\varphi_n)[s]\) iff \(p(e_1,\ldots,e_n)=1\), where \(e_i=1\) if
\(\mathfrak M\models\varphi_i[s]\) and \(e_i=0\) otherwise.

\begin{hbtheorem}{thm:p02-w02c-p02-003-or-not}{proved=annex:p02-w02c-p02-003-or-not}{Every finitary connective of positive arity is definable from \(\lor\) and \(\neg\)}
For every \(n\geq1\) and every \(p:\{0,1\}^{n}\to\{0,1\}\) there is a
propositional formula \(\pi\) built from \(p_1,\ldots,p_n\) with \(\lor\)
and \(\neg\) alone whose truth function is \(p\).  Consequently, for all
first-order formulas \(\varphi_1,\ldots,\varphi_n\) of a language \(L\),
the formula \(\pi(\varphi_1,\ldots,\varphi_n)\) obtained by substituting
\(\varphi_i\) for \(p_i\) satisfies, for every \(L\)-structure
\(\mathfrak M\) and assignment \(s\),
\(\mathfrak M\models\pi(\varphi_1,\ldots,\varphi_n)[s]\) iff
\(\mathfrak M\models\#(\varphi_1,\ldots,\varphi_n)[s]\).
\end{hbtheorem}

A nullary connective (a truth constant) has no arguments to substitute, so
no formula built from \(\lor\) and \(\neg\) alone represents it; this is
\cref{thm:p02-w02-p02-008-nullary}.  In first-order logic the two constants
are expressed by sentences such as \(\forall x\,x=x\) and its negation, so
the remark holds for them with the first-order quantifiers and identity
added.

Barwise also remarks that the natural structure for the language of set
theory, whose domain is the collection of all sets, is not itself a set,
and that by the Completeness Theorem a set of axioms with a model in any
reasonable sense has a reasonably small set-theoretic model
\cite[ch.~A.1, \S3, pp.~17--18]{Barwise1977FirstOrder}.  The following
theorem makes ``reasonable'' precise as soundness of the proof calculus.

\begin{hbtheorem}{thm:p02-w02c-p02-003-small-model}{proved=annex:p02-w02c-p02-003-small-model}{Small set-theoretic models (Barwise, p.~18)}
Assume the axiom of choice.  Let \(L\) be a language and \(T\) a set of
first-order sentences of \(L\).  Suppose there is some notion of truth for
\(L\)-sentences, in an object \(\mathcal M\) of any kind, such that every
sentence derivable from \(T\) in the natural deduction calculus of
\cref{thm:fol-completeness} is true in \(\mathcal M\), while the sentence
\(\neg\forall x\,x=x\) is not.  Then \(T\) has a set-theoretic model whose
universe has cardinality at most \(\max(\lvert L\rvert,\aleph_0)\).
\end{hbtheorem}

\begin{hbwarning}{def:p02-w02c-p02-003-small-model-limit}{The hypothesis is about all derivations at once}
The hypothesis of \cref{thm:p02-w02c-p02-003-small-model} is one statement
about every derivation from \(T\).  For Barwise's example, the axioms of
set theory and the collection of all sets with membership, truth in that
collection is expressed one formula at a time, so the hypothesis is not
available as a single statement about it, and the theorem does not yield a
set-theoretic model in that case.  Barwise's remark is established here in
the form of the theorem; the case of a class-sized structure for which only
the axioms are known to be true, one at a time, is not.
\end{hbwarning}

\subsection{Barwise's further claims about logics beyond first order}
\label{sec:p02-w02c-p02-003-beyond}

Barwise reports, at \S5, that the natural numbers and finite sequences are
definable in \(\mathrm{HF}(\mathfrak M)\) and that
\(\mathrm{HF}(\mathfrak M)\) is admissible (p.~43); that a completeness
theorem for countable theories of \(L_{\omega_1\omega}\) needs an
infinitary rule; that weak second-order logic translates easily into
\(L_{\omega_1\omega}\) but not conversely; that weak second-order logic has
relations that are implicitly but not explicitly definable (p.~44); that
\(L(Q_0)\) is essentially equivalent to \(\omega\)-logic and weak
second-order logic (p.~44); and that \(L(Q)\) has a Completeness and a
Compactness Theorem for finite or countable languages (pp.~44--45)
\cite[ch.~A.1, \S5]{Barwise1977FirstOrder}.

\begin{hbconvention}{def:p02-w02c-p02-003-wso-reading}{Atomic formulas of \(L^{*}\) and admissibility}
The chapter does not list the atomic formulas of \(L^{*}\).  They are read
here as the atomic formulas of \(L\), and \(t\in a\), \(a\in b\) and
\(a=b\) for an \(L\)-term \(t\) and set variables \(a,b\), with the
quantifiers \(\forall a\), \(\exists a\) over \(\mathrm{HF}(M)\); the
elements of \(M\) are atoms, as in \cref{def:p02-w02-p02-005-bw-weak-so}.
For \(X\subseteq M\), \(\mathrm{HF}(X)\) is defined as
\(\mathrm{HF}(M)\) with \(X\) in place of \(M\).  Admissibility is that of
Makkai's chapter, to which Barwise refers: an \emph{admissible set} is a
transitive set \(A\), in the universe of sets over a set \(U\) of
urelements, that is a model of the Kripke--Platek axioms with urelements
\(\mathrm{KPU}\) in the language \(\{\in,U\}\): urelements have no members,
empty set, extensionality, foundation, pair, union,
\(\Delta_0\)-separation and \(\Delta_0\)-collection, where a
\(\Delta_0\)-formula is one whose quantifiers are all bounded, of the form
\(\forall x\in y\) or \(\exists x\in y\)
\cite[ch.~A.7, \S2, 2.1, 2.3, 2.4, pp.~238--239]{Makkai1977}.
\end{hbconvention}

\begin{hbtheorem}{thm:p02-w02c-p02-003-hf}{proved=annex:p02-w02c-p02-003-hf}{Natural numbers, finite sequences and admissibility in \(\mathrm{HF}(\mathfrak M)\)}
Let \(\mathfrak M\) be a structure for \(L\) with domain \(M\), and let
\(\mathrm{Nat}(a)\) be the \(L^{*}\)-formula
\[
\begin{gathered}
 \neg\exists x\,(x\in a)\land
 \forall b\in a\,\forall x\,(x\in b\to x\in a)\\
 {}\land\forall b\in a\,\forall c\in b\,(c\in a)\land
 \forall b\in a\,\forall c\in b\,\forall d\in c\,(d\in b).
\end{gathered}
\]
\begin{enumerate}
\item Every von Neumann natural number \(n=\{0,\ldots,n-1\}\) lies in
      \(\mathrm{HF}(M)\), and for \(a\in\mathrm{HF}(M)\),
      \(\mathrm{HF}(\mathfrak M)\models\mathrm{Nat}[a]\) iff \(a\) is a von
      Neumann natural number.
\item For every \(n\) and every \(v_0,\ldots,v_{n-1}\in M\cup\mathrm{HF}(M)\),
      the set \(\{\langle i,v_i\rangle:i<n\}\) of Kuratowski pairs
      \(\langle i,v\rangle=\{\{i\},\{i,v\}\}\) lies in \(\mathrm{HF}(M)\).
      There are \(L^{*}\)-formulas \(\mathrm{Seq}(s,n)\),
      \(\mathrm{Val}(s,i,x)\) and \(\mathrm{Val}'(s,i,b)\) such that, for
      \(s,n,i,b\in\mathrm{HF}(M)\) and \(x\in M\),
      \(\mathrm{Seq}(s,n)\) holds in \(\mathrm{HF}(\mathfrak M)\) iff \(n\)
      is a natural number and \(s\) is a function with domain \(n\); and,
      when \(\mathrm{Seq}(s,n)\) holds, \(\mathrm{Val}(s,i,x)\) holds iff
      \(i\in n\) and \(s(i)=x\), and \(\mathrm{Val}'(s,i,b)\) holds iff
      \(i\in n\) and \(s(i)=b\).
\item The set \(A=M\cup\mathrm{HF}(M)\), with \(U=M\) and the membership
      relation, is an admissible set.  Moreover every instance of
      separation and collection, for arbitrary formulas of \(\{\in,U\}\)
      with relation symbols for the relations of \(\mathfrak M\) added,
      holds in \(A\).
\end{enumerate}
\end{hbtheorem}

Barwise adds that, \(\mathrm{HF}(\mathfrak M)\) being admissible, functions
can be defined by recursion, and that the sentences of his \S2 called weak
second-order are weak second-order in the precise sense of \(L^{*}\)
\cite[ch.~A.1, \S5, 5.3, p.~43]{Barwise1977FirstOrder}; he refers to
Makkai's recursion principles \cite[ch.~A.7, \S2, 2.16--2.17, p.~243]{Makkai1977}.
The two results below prove the form of recursion that these sentences
use, recursion along the natural numbers, directly from parts~1 and~2 of
\cref{thm:p02-w02c-p02-003-hf}, and then write each of the \S2 sentences
as an \(L^{*}\)-sentence.  Makkai's \(\Sigma\)-recursion on membership is
proved for \(\mathrm{HF}(\mathfrak M)\) afterwards, in
\cref{thm:p02-w02c-p02-003-hf-sigma-recursion}.
In the formulas below, \(\forall i\in n\,\chi\) abbreviates
\(\forall i\,(i\in n\to\chi)\), and
\begin{align*}
 \mathrm{Emp}(e)\colon\;&\neg\exists y\,(y\in e)\land\neg\exists b\,(b\in e),\\
 \mathrm{Succ}(i,j)\colon\;&\forall y\,(y\in j\leftrightarrow y\in i)\land
   \forall b\,(b\in j\leftrightarrow(b\in i\lor b=i)),
\end{align*}
with \(y\) an individual and \(b,e,i,j\) set variables; in
\(\mathrm{HF}(\mathfrak M)\), \(\mathrm{Emp}(e)\) says \(e=\emptyset\) and
\(\mathrm{Succ}(i,j)\) says \(j=i\cup\{i\}\).

\begin{hbtheorem}{thm:p02-w02c-p02-003-hf-recursion}{proved=annex:p02-w02c-p02-003-hf-recursion}{Definition by recursion on the natural numbers in \(\mathrm{HF}(\mathfrak M)\)}
Let \(\mathfrak M\) be a structure for \(L\) with domain \(M\).  Fix a
\emph{value sort}: either the individuals (values in \(M\)) or the sets
(values in \(\mathrm{HF}(M)\)); write \(\mathrm{Val}^{\ast}\) for
\(\mathrm{Val}\) in the first case and \(\mathrm{Val}'\) in the second
(part~2 of \cref{thm:p02-w02c-p02-003-hf}).  Let
\(\gamma(i,v,w,\bar p)\) be an \(L^{*}\)-formula, with \(i\) a set
variable, \(v,w\) variables of the value sort and parameters \(\bar p\),
and fix \(\bar p\) in \(M\cup\mathrm{HF}(M)\) such that for every natural
number \(i\) and every \(v\) of the value sort there is exactly one \(w\)
of the value sort with \(\mathrm{HF}(\mathfrak M)\models\gamma[i,v,w,\bar
p]\); call it \(G(i,v)\).  For \(v_0\) of the value sort, let
\(F_{v_0}(0)=v_0\) and \(F_{v_0}(i+1)=G(i,F_{v_0}(i))\).  Let
\(\mathrm{Rec}_\gamma(n,v_0,w,\bar p)\) be
\[
\begin{aligned}
&\mathrm{Nat}(n)\land\exists m\,\exists s\,\bigl[\mathrm{Succ}(n,m)\land
 \mathrm{Seq}(s,m)\land\mathrm{Val}^{\ast}(s,n,w)\\
&\quad{}\land\exists e\,(\mathrm{Emp}(e)\land\mathrm{Val}^{\ast}(s,e,v_0))
 \land\forall i\in n\,\forall j\in m\,\bigl(\mathrm{Succ}(i,j)\to
 \forall v\,\forall w'\,\delta\bigr)\bigr],
\end{aligned}
\]
where \(\delta\) is \((\mathrm{Val}^{\ast}(s,i,v)\land
\mathrm{Val}^{\ast}(s,j,w'))\to\gamma(i,v,w',\bar p)\).
Then for every \(n\in\mathrm{HF}(M)\) and \(v_0,w\) of the value sort,
\(\mathrm{HF}(\mathfrak M)\models\mathrm{Rec}_\gamma[n,v_0,w,\bar p]\) iff
\(n\) is a natural number and \(w=F_{v_0}(n)\).
\end{hbtheorem}

For \(n\geq1\) and a term \(x\), Barwise's term \(nx\) is \(x\) for
\(n=1\) and \((mx+x)\) for \(n=m+1\); the term \(n1\) is the case
\(x=1\) \cite[ch.~A.1, \S2, pp.~8, 12]{Barwise1977FirstOrder}.  For a
structure in which \(+\) is interpreted, \((nx)^{\mathfrak M}\) is its
value.  Barwise's sentences are
\begin{align*}
&\forall n\geq1\,\forall x\,\exists y\,[ny=x] &&\text{(6), divisible,}\\
&\forall n\geq1\,\forall x\,[x\neq0\to nx\neq0] &&\text{(7), torsion-free,}\\
&\forall x\,\exists n\geq1\,[nx=0] &&\text{(8), torsion,}\\
&\forall p\,[p\text{ a prime}\to p1\neq0] &&\text{(10), characteristic }0,
\end{align*}
with \(n\) and \(p\) ranging over the natural numbers
\cite[ch.~A.1, \S2, pp.~8--9, 15]{Barwise1977FirstOrder}.  The
\(L^{*}\)-formulas used to write them are built with
\cref{thm:p02-w02c-p02-003-hf-recursion}:
\begin{itemize}
\item \(\mathrm{Mul}(n,x,y)\) is \(\exists k\,(\mathrm{Succ}(k,n)\land
      \mathrm{Rec}_{\gamma_+}(k,x,y,x))\), where the value sort is the
      individuals and \(\gamma_+(i,v,w,x)\) is \(w=v+x\);
\item \(\mathrm{Add}(a,i,w)\) is the formula
      \(\mathrm{Rec}_{\sigma}(i,a,w)\), where the value sort is the sets and \(\sigma(i,v,w)\) is
      \(\mathrm{Succ}(v,w)\);
\item \(\mathrm{Times}(a,b,w)\) is \(\exists e\,(\mathrm{Emp}(e)\land
      \mathrm{Rec}_{\tau}(b,e,w,a))\), where the value sort is the sets and
      \(\tau(i,v,w,a)\) is \(\mathrm{Add}(v,a,w)\);
\item \(\mathrm{Prime}(p)\) is \(\mathrm{Nat}(p)\land\exists b\in p\,
      \exists c\,(c\in b)\land\forall a\in p\,\forall b\in p\,
      \neg\mathrm{Times}(a,b,p)\);
\item \(\mathrm{Pos}(n)\) is \(\mathrm{Nat}(n)\land\exists b\,(b\in n)\).
\end{itemize}

\begin{hbtheorem}{thm:p02-w02c-p02-003-wso-sentences}{proved=annex:p02-w02c-p02-003-wso-sentences}{The \S2 sentences are weak second-order (Barwise, p.~43)}
Let \(\mathfrak M\) be a structure for a language \(L\) containing \(+\)
and \(0\), and let
\begin{align*}
(6)^{*}\colon\;&\forall n\,(\mathrm{Pos}(n)\to\forall x\,\exists y\,
  \mathrm{Mul}(n,y,x)),\\
(7)^{*}\colon\;&\forall n\,(\mathrm{Pos}(n)\to\forall x\,\forall y\,
  ((\neg x=0\land\mathrm{Mul}(n,x,y))\to\neg y=0)),\\
(8)^{*}\colon\;&\forall x\,\exists n\,(\mathrm{Pos}(n)\land
  \mathrm{Mul}(n,x,0)),\\
(10)^{*}\colon\;&\forall p\,(\mathrm{Prime}(p)\to\forall y\,
  (\mathrm{Mul}(p,1,y)\to\neg y=0)),
\end{align*}
the last when \(L\) also contains \(1\).  Then
\(\mathrm{HF}(\mathfrak M)\models(6)^{*}\) iff for every natural number
\(n\geq1\) and every \(x\in M\) there is \(y\in M\) with
\((ny)^{\mathfrak M}=x\); \(\mathrm{HF}(\mathfrak M)\models(7)^{*}\) iff
\((nx)^{\mathfrak M}\neq0^{\mathfrak M}\) for every \(n\geq1\) and
\(x\neq0^{\mathfrak M}\); \(\mathrm{HF}(\mathfrak M)\models(8)^{*}\) iff
every \(x\in M\) has some \(n\geq1\) with
\((nx)^{\mathfrak M}=0^{\mathfrak M}\); and
\(\mathrm{HF}(\mathfrak M)\models(10)^{*}\) iff
\((p1)^{\mathfrak M}\neq0^{\mathfrak M}\) for every prime number \(p\).
\end{hbtheorem}

The Archimedean axiom, the remaining sentence that Barwise calls weak
second-order (p.~12), is written as an \(L^{*}\)-sentence
\(\mathrm{Arch}^{*}\) in the proof of
\cref{thm:p02-w02-p02-005-bw-countable-field}.  Barwise replaces
\((10)\) by the infinite list \((10)_p\colon p1\neq0\), one first-order
sentence for each prime \(p\) \cite[ch.~A.1, \S2, p.~15]{Barwise1977FirstOrder};
by the theorem, \(\mathrm{HF}(\mathfrak M)\models(10)^{*}\) iff
\(\mathfrak M\) satisfies every \((10)_p\).

The recursion of \cref{thm:p02-w02c-p02-003-hf-recursion} runs along the
natural numbers.  Makkai's principle runs along membership: the value at
\(y\) is computed from the values at all members of the transitive
closure of \(y\) \cite[ch.~A.7, \S2, 2.16--2.17, p.~243]{Makkai1977}.  The
next theorem proves it for the admissible set \(A=M\cup\mathrm{HF}(M)\) of
part~3 of \cref{thm:p02-w02c-p02-003-hf}, regarded as the structure
\((A,\in,U)\) with \(U=M\), possibly with relation symbols for relations
of \(\mathfrak M\) added.  A set \(b\) is \emph{transitive} if every member
of \(b\) that is a set has all its members in \(b\).  For a set \(a\in
A\), \(\mathrm{TC}(a)\) is the least transitive set \(b\) with
\(a\subseteq b\), and \(\mathrm{TC}(p)=\emptyset\) for an urelement
\(p\in M\) \cite[ch.~A.7, \S2, p.~242]{Makkai1977}.  A
\(\Sigma\)-formula is one built from \(\Delta_0\)-formulas
(\cref{def:p02-w02c-p02-003-wso-reading}) by conjunction, disjunction,
bounded quantification and existential quantification
\cite[ch.~A.7, \S2, 2.5, p.~239]{Makkai1977}; an operation on \(A\) is
\(\Sigma\) if its graph is defined in \((A,\in,U)\) by a
\(\Sigma\)-formula with parameters in \(A\)
\cite[ch.~A.7, \S2, 2.8, p.~240]{Makkai1977}.  The following
\(\Delta_0\)-formulas of \(\{\in,U\}\) are used, with all bound variables
new:
\begin{itemize}
\item \(\mathrm{Pr}(p,z,v)\), saying \(p=\{\{z\},\{z,v\}\}\):
      \(\exists q\in p\,\exists q'\in p\,[z\in q\land\forall u\in q\,(u=z)
      \land z\in q'\land v\in q'\land\forall u\in q'\,(u=z\lor u=v)\land
      \forall r\in p\,(r=q\lor r=q')]\);
\item \(\mathrm{Ap}(f,z,v)\), saying \(\langle z,v\rangle\in f\):
      \(\exists p\in f\,\mathrm{Pr}(p,z,v)\); and
      \(\exists^{f}v\,(z,v)\,\chi\) abbreviates
      \(\exists p\in f\,\exists q\in p\,\exists v\in q\,
      (\mathrm{Pr}(p,z,v)\land\chi)\);
\item \(\mathrm{Fun}(f,t)\), saying that \(f\) is a function with domain
      \(t\): the conjunction of \(\neg U(f)\), \(\forall p\in f\,\exists
      z\in t\,\exists q\in p\,\exists v\in q\,\mathrm{Pr}(p,z,v)\),
      \(\forall z\in t\,\exists^{f}v\,(z,v)\,(v=v)\) and
      \(\forall z\in t\,\forall p\in f\,\forall p'\in f\,\forall q\in p\,
      \forall q'\in p'\,\forall v\in q\,\forall v'\in q'\,
      ((\mathrm{Pr}(p,z,v)\land\mathrm{Pr}(p',z,v'))\to v=v')\);
\item \(\mathrm{Tr}(t)\): \(\neg U(t)\land\forall u\in t\,\forall v\in u\,
      (v\in t)\);
\item \(\mathrm{Cl}(h,t)\), saying that \(h(z)\) is the set
      \(z\cup\bigcup_{u\in z}h(u)\) for every \(z\in t\):
      \(\forall z\in t\,\exists^{h}c\,(z,c)\,[\neg U(c)\land
      \forall v\in z\,(v\in c)\land\forall u\in z\,\exists^{h}c'\,(u,c')\,
      \forall v\in c'\,(v\in c)\land
      \forall v\in c\,(v\in z\lor\exists u\in z\,\exists^{h}c'\,(u,c')\,
      (v\in c'))]\);
\item \(\mathrm{Rs}(H,f,h,t)\), saying that \(H(z)=f\restriction h(z)\) for
      every \(z\in t\):
      \(\forall z\in t\,\exists^{H}g\,(z,g)\,\exists^{h}c\,(z,c)\,
      [\neg U(g)\land\forall p\in g\,(p\in f\land\exists u\in c\,\exists
      q\in p\,\exists v\in q\,\mathrm{Pr}(p,u,v))\land
      \forall u\in c\,\forall p\in f\,\forall q\in p\,\forall v\in q\,
      (\mathrm{Pr}(p,u,v)\to p\in g)]\).
\end{itemize}

\begin{hbtheorem}{thm:p02-w02c-p02-003-hf-sigma-recursion}{proved=annex:p02-w02c-p02-003-hf-sigma-recursion}{\(\Sigma\)-recursion on membership in \(\mathrm{HF}(\mathfrak M)\) (Makkai, 2.16--2.17)}
Let \(\mathfrak M\) be a structure with domain \(M\), let
\(A=M\cup\mathrm{HF}(M)\), let \(n\geq0\) and let \(G\colon
A^{n+2}\to A\).  For \(F\colon A^{n+1}\to A\), \(\bar x\in A^{n}\) and
\(y\in A\) write \(F_{\bar x}\restriction\mathrm{TC}(y)\) for the set
\(\{\langle z,F(\bar x,z)\rangle:z\in\mathrm{TC}(y)\}\).
\begin{enumerate}
\item For \(y\in A\), \(\mathrm{TC}(y)\) exists, is a finite subset of
      \(A\) and is a member of \(\mathrm{HF}(M)\), and so is \(F_{\bar
      x}\restriction\mathrm{TC}(y)\) for every \(F\).  There is exactly one
      \(F\colon A^{n+1}\to A\) such that, for all \(\bar x\) and \(y\),
      \[
       F(\bar x,y)=G(\bar x,y,F_{\bar x}\restriction\mathrm{TC}(y)).
      \]
\item Suppose that \(\gamma(\bar x,y,g,w,\bar r)\) is a formula of
      \(\{\in,U\}\), possibly with relation symbols for relations of
      \(\mathfrak M\), and \(\bar r\in A\) are parameters such that
      \(G(\bar x,y,g)=w\) iff \((A,\in,U,\ldots)\models\gamma[\bar
      x,y,g,w,\bar r]\).  Let \(\Phi_\gamma(\bar x,y,w,\bar r)\) be
      \[
      \begin{aligned}
      &\exists t\,\exists f\,\exists h\,\exists H\,\bigl[\mathrm{Tr}(t)\land
       y\in t\land\mathrm{Fun}(f,t)\land\mathrm{Fun}(h,t)\\
      &\qquad{}\land\mathrm{Fun}(H,t)
       \land\mathrm{Cl}(h,t)\land\mathrm{Rs}(H,f,h,t)\\
      &\qquad{}\land\forall z\in t\,\exists^{H}g\,(z,g)\,
       \exists^{f}v\,(z,v)\,\gamma(\bar x,z,g,v,\bar r)
       \land\mathrm{Ap}(f,y,w)\bigr].
      \end{aligned}
      \]
      Then \(F(\bar x,y)=w\) iff \((A,\in,U,\ldots)\models\Phi_\gamma[\bar
      x,y,w,\bar r]\).
\item If \(\gamma\) is a \(\Sigma\)-formula, so is \(\Phi_\gamma\); hence
      if \(G\) is a \(\Sigma\) operation on the admissible set \(A\), so is
      \(F\).
\item Every formula \(\varphi(v_1,\ldots,v_k)\) of \(\{\in,U\}\) (with
      relation symbols for relations of \(\mathfrak M\)) has, for each
      choice of a sort (individual or set) for each \(v_i\), an
      \(L^{*}\)-formula \(\varphi^{\mathrm s}\) with
      \((A,\in,U,\ldots)\models\varphi[\bar a]\) iff
      \(\mathrm{HF}(\mathfrak M)\models\varphi^{\mathrm s}[\bar a]\) for all
      \(\bar a\) of the chosen sorts.  So the graph of \(F\), restricted to
      arguments and values of fixed sorts, is defined in
      \(\mathrm{HF}(\mathfrak M)\) by an \(L^{*}\)-formula.
\end{enumerate}
\end{hbtheorem}

Part~1 is the identity of Makkai's 2.16, and part~3 is the conclusion of
his 2.17, for this admissible set.  Makkai states 2.16 as a theorem of the
axiom system \(\mathrm{KPU}\) and refers its proof to the recursion
principle of set theory.  The proof in the annex uses only the finiteness
of the sets in \(\mathrm{HF}(M)\) and induction on the stage at which a set
enters \(\mathrm{HF}(M)\), so it does not pass through \(\mathrm{KPU}\).

\begin{hbtheorem}{thm:p02-w02c-p02-003-infinitary-rule}{proved=annex:p02-w02c-p02-003-infinitary-rule}{A complete proof system for \(L_{\omega_1\omega}\) needs a rule with infinitely many premises}
Call a relation \(\vdash\) between sets of sentences of
\(L_{\omega_1\omega}\) and such sentences \emph{finitary} if
\(\Gamma\vdash\varphi\) implies \(\Gamma_0\vdash\varphi\) for some finite
\(\Gamma_0\subseteq\Gamma\), and \emph{sound} if \(\Gamma\vdash\varphi\)
implies that every model of \(\Gamma\) satisfies \(\varphi\).  No sound
finitary relation is complete for countable theories: for the empty
language there are a countable set \(T\) of sentences and a sentence
\(\bot\) such that every model of \(T\) satisfies \(\bot\) and
\(T\nvdash\bot\).  In particular this holds for derivability in any proof
system in which each derivation is a finite object using finitely many
premises.
\end{hbtheorem}

This is the necessity half of Barwise's remark.  A calculus with a rule of
countably many premises that is complete for countable theories is
\cref{def:p02-w02c-p02-003-hinf}, with completeness
\cref{thm:p02-w02c-p02-003-infinitary-completeness}, after the model
existence theorem below.

\begin{hbtheorem}{thm:p02-w02c-p02-003-wso-translation}{proved=annex:p02-w02c-p02-003-wso-translation}{Weak second-order logic translates into \(L_{\omega_1\omega}\), and not conversely}
\begin{enumerate}
\item For every language \(L\) and every sentence \(\varphi\) of
      \(L^{*}\) there is a sentence \(\varphi^{\natural}\) of
      \(L_{\omega_1\omega}\) for \(L\) such that, for every structure
      \(\mathfrak M\) for \(L\), \(\mathrm{HF}(\mathfrak M)\models\varphi\)
      iff \(\mathfrak M\models\varphi^{\natural}\).
\item For the empty language, with countably many variables of each
      sort, there is a sentence \(\psi\) of \(L_{\omega_1\omega}\) such that
      no sentence \(\varphi\) of \(L^{*}\) satisfies
      \(\mathrm{HF}(\mathfrak M)\models\varphi\) iff
      \(\mathfrak M\models\psi\) for every \(\mathfrak M\).
\end{enumerate}
\end{hbtheorem}

Barwise contrasts weak second-order logic, where some relations are
implicitly but not explicitly definable, with first-order logic and
\(L_{\omega_1\omega}\), where by Beth's theorem this cannot happen
\cite[ch.~A.1, \S5, p.~44]{Barwise1977FirstOrder}.  For first-order logic
this is \cref{thm:p03-beth-definability}.  For \(L_{\omega_1\omega}\)
Makkai states Beth's theorem for single sentences and attributes it to
Lopez-Escobar \cite[ch.~A.7, \S1, p.~236]{Makkai1977}, and derives it from
the interpolation theorem \cite[ch.~A.7, \S8, 8.6, p.~274]{Makkai1977}.
Makkai proves interpolation from his game-form theorem and Vaught's
covering theorem \cite[ch.~A.7, \S8, 8.1--8.2, pp.~271--272]{Makkai1977}
and names a second route through consistency properties
\cite[ch.~A.7, \S8, p.~274]{Makkai1977}.  The proof below takes the second
route.  Its models are built from Hintikka sets in Makkai's sense
\cite[ch.~A.7, \S4, 4.1--4.2, p.~250]{Makkai1977}, and its model existence
theorem plays the part of Makkai's 4.3
\cite[ch.~A.7, \S4, p.~251]{Makkai1977}.  A sentence \(\chi\) of
\(L_{\omega_1\omega}\) for a language \(K\) is \emph{valid} if every
structure for \(K\) satisfies it, \(\chi\to\chi'\) abbreviates
\(\neg\chi\lor\chi'\), and the finitary \(\land\), \(\lor\) are read as
\(\bigwedge\), \(\bigvee\) of two-element sets.  For a set \(\Gamma\) of
sentences and a sentence \(\theta\), all for a language \(K\),
\(\Gamma\models\theta\) means that every structure for \(K\) satisfying
every member of \(\Gamma\) satisfies \(\theta\).  For a formula \(\chi\), a
variable \(x\) and a constant symbol \(c\), \(\chi(c/x)\) is the result of
replacing every free occurrence of \(x\) in \(\chi\) by \(c\).

\begin{hbdefinition}{def:p02-w02c-p02-003-hintikka}{Negation normal form and Hintikka sets (Makkai, 4.1)}
A formula of \(L_{\omega_1\omega}\) is in \emph{negation normal form}
(n.n.f.) if it is built from atomic and negated atomic formulas by
\(\bigwedge\), \(\bigvee\), \(\forall\) and \(\exists\) only
\cite[ch.~A.7, \S4, p.~249]{Makkai1977}.  For a formula \(\chi\), the
n.n.f. formulas \(\chi^{+}\) and \(\chi^{\ominus}\) are defined together
by recursion on \(\chi\): for atomic \(\chi\), \(\chi^{+}=\chi\) and
\(\chi^{\ominus}=\neg\chi\); \((\neg\chi)^{+}=\chi^{\ominus}\) and
\((\neg\chi)^{\ominus}=\chi^{+}\);
\((\bigwedge\Phi)^{+}=\bigwedge\{\varphi^{+}:\varphi\in\Phi\}\) and
\((\bigwedge\Phi)^{\ominus}=\bigvee\{\varphi^{\ominus}:\varphi\in\Phi\}\);
\((\bigvee\Phi)^{+}=\bigvee\{\varphi^{+}:\varphi\in\Phi\}\) and
\((\bigvee\Phi)^{\ominus}=\bigwedge\{\varphi^{\ominus}:\varphi\in\Phi\}\);
\((\exists x\,\chi)^{+}=\exists x\,\chi^{+}\),
\((\exists x\,\chi)^{\ominus}=\forall x\,\chi^{\ominus}\),
\((\forall x\,\chi)^{+}=\forall x\,\chi^{+}\) and
\((\forall x\,\chi)^{\ominus}=\exists x\,\chi^{\ominus}\).

Let \(K\) be a relational language and \(C\) a nonempty set of constant
symbols not in \(K\); the closed terms of \(K\cup C\) are the members of
\(C\).  A set \(H\) of n.n.f. sentences of \(L_{\omega_1\omega}\) for
\(K\cup C\) is a \emph{Hintikka set} if
\begin{enumerate}
\item[(i)] for no atomic sentence \(\theta\) are both \(\theta\) and
      \(\neg\theta\) in \(H\);
\item[(iia)] \((c=c)\in H\) for every \(c\in C\);
\item[(iib)] if \(\varphi(v)\) is an atomic formula with no variable other
      than \(v\), \(\varphi(c/v)\in H\) and \((c=d)\in H\), then
      \(\varphi(d/v)\in H\);
\item[(iii)] if \(\bigwedge\Phi\in H\) then \(\Phi\subseteq H\);
\item[(iv)] if \(\forall x\,\varphi\in H\) then \(\varphi(c/x)\in H\) for
      every \(c\in C\);
\item[(v)] if \(\bigvee\Phi\in H\) then \(\varphi\in H\) for some
      \(\varphi\in\Phi\);
\item[(vi)] if \(\exists x\,\varphi\in H\) then \(\varphi(c/x)\in H\) for
      some \(c\in C\).
\end{enumerate}
These are Makkai's conditions 4.1 (i)--(vi)
\cite[ch.~A.7, \S4, 4.1, p.~250]{Makkai1977} for a language whose closed
terms are its constants.
\end{hbdefinition}

\begin{hblemma}{thm:p02-w02c-p02-003-nnf}{proved=annex:p02-w02c-p02-003-nnf}{Negation normal form}
For every formula \(\chi\) of \(L_{\omega_1\omega}\), \(\chi^{+}\) and
\(\chi^{\ominus}\) are n.n.f. formulas of \(L_{\omega_1\omega}\) with the
same free variables and the same non-logical symbols as \(\chi\), and for
every structure \(\mathfrak M\) and assignment \(s\),
\(\mathfrak M\models\chi^{+}[s]\) iff \(\mathfrak M\models\chi[s]\), and
\(\mathfrak M\models\chi^{\ominus}[s]\) iff
\(\mathfrak M\not\models\chi[s]\).
\end{hblemma}

\begin{hblemma}{thm:p02-w02c-p02-003-hintikka-model}{proved=annex:p02-w02c-p02-003-hintikka-model}{A Hintikka set has a model (Makkai, 4.2)}
Let \(K\) be a relational language, \(C\) a nonempty set of constant
symbols not in \(K\), and \(H\) a Hintikka set for \(K\cup C\).  There is a
structure \(\mathfrak M\) for \(K\cup C\) that satisfies every member of
\(H\) and in which every element is \(c^{\mathfrak M}\) for some \(c\in C\).
Makkai states this for arbitrary languages and omits the induction steps of
its proof \cite[ch.~A.7, \S4, 4.2, p.~250]{Makkai1977}.
\end{hblemma}

\begin{hbdefinition}{def:p02-w02c-p02-003-consistency}{Consistency property}
Let \(K\) be a relational language and \(C\) a countably infinite set of
constant symbols not in \(K\).  A \emph{consistency property} for
\(K\cup C\) is a set \(\mathcal S\) of finite sets of n.n.f. sentences of
\(L_{\omega_1\omega}\) for \(K\cup C\) such that for every
\(s\in\mathcal S\):
\begin{enumerate}
\item[(C0)] for no atomic sentence \(\theta\) are both \(\theta\) and
      \(\neg\theta\) in \(s\);
\item[(C1)] \(s\cup\{c=c\}\in\mathcal S\) for every \(c\in C\);
\item[(C2)] if \(\varphi(v)\) is an atomic formula with no variable other
      than \(v\), and \(\varphi(c/v)\in s\) and \((c=d)\in s\), then
      \(s\cup\{\varphi(d/v)\}\in\mathcal S\);
\item[(C3)] if \(\bigwedge\Phi\in s\) and \(\varphi\in\Phi\), then
      \(s\cup\{\varphi\}\in\mathcal S\);
\item[(C4)] if \(\forall x\,\varphi\in s\) and \(c\in C\), then
      \(s\cup\{\varphi(c/x)\}\in\mathcal S\);
\item[(C5)] if \(\bigvee\Phi\in s\), then
      \(s\cup\{\varphi\}\in\mathcal S\) for some \(\varphi\in\Phi\);
\item[(C6)] if \(\exists x\,\varphi\in s\), then
      \(s\cup\{\varphi(c/x)\}\in\mathcal S\) for some \(c\in C\).
\end{enumerate}
\end{hbdefinition}

\begin{hbtheorem}{thm:p02-w02c-p02-003-model-existence}{proved=annex:p02-w02c-p02-003-mex}{Model existence}
Assume the axiom of choice.  Let \(K\) be a countable relational language,
\(C\) a countably infinite set of constant symbols not in \(K\), and
\(\mathcal S\) a consistency property for \(K\cup C\).  Then every
\(s\in\mathcal S\) has a model: a structure for \(K\cup C\) that satisfies
every member of \(s\).
\end{hbtheorem}

\begin{hbtheorem}{thm:p02-w02c-p02-003-interpolation}{proved=annex:p02-w02c-p02-003-interpolation}{Lopez-Escobar interpolation for \(L_{\omega_1\omega}\)}
Assume the axiom of choice.  Let \(L\) be a language in Barwise's sense
(relation, function and constant symbols), and let \(\varphi\) and
\(\psi\) be sentences of \(L_{\omega_1\omega}\) for \(L\) such that
\(\varphi\to\psi\) is valid.  Then there is a sentence \(\theta\) of
\(L_{\omega_1\omega}\), a \emph{Craig interpolant} of \(\varphi\) and
\(\psi\), such that \(\varphi\to\theta\) and \(\theta\to\psi\) are both
valid and every non-logical symbol of \(\theta\) occurs both in
\(\varphi\) and in \(\psi\).  This is the first assertion of Makkai's
8.2(i) \cite[ch.~A.7, \S8, 8.2(i), p.~272]{Makkai1977}.
\end{hbtheorem}

\begin{hbtheorem}{thm:p02-w02c-p02-003-beth}{proved=annex:p02-w02c-p02-003-beth}{Beth's definability theorem for \(L_{\omega_1\omega}\)}
Let \(L\) be a language, \(P\) an \(n\)-ary relation symbol not in \(L\),
and \(\varphi(P)\) a sentence of \(L_{\omega_1\omega}\) for
\(L\cup\{P\}\).  Suppose that \(\varphi(P)\) defines \(P\)
\emph{implicitly}: for every structure \(\mathfrak M\) for \(L\) there is
at most one relation \(R\subseteq M^{n}\) with
\((\mathfrak M,R)\models\varphi(P)\).  Let the set of variables be
infinite, and let \(x_1,\ldots,x_n\) be distinct variables.  Then \(P\) is
\emph{explicitly} definable: there is a formula \(\theta(x_1,\ldots,x_n)\)
of \(L_{\omega_1\omega}\), with free variables among \(x_1,\ldots,x_n\) and
with every non-logical symbol occurring in \(\varphi(P)\) and belonging to
\(L\), such that every model of \(\varphi(P)\) satisfies
\[
 \forall x_1\cdots\forall x_n\,
 \bigl(P(x_1,\ldots,x_n)\leftrightarrow\theta(x_1,\ldots,x_n)\bigr).
\]
The same holds with a countable set \(\Phi(P)\) of sentences of
\(L_{\omega_1\omega}\) for \(L\cup\{P\}\) in place of \(\varphi(P)\).  The
axiom of choice is assumed; it enters through
\cref{thm:p02-w02c-p02-003-interpolation}.
\end{hbtheorem}

\begin{hbexample}{def:p02-w02c-p02-003-beth-example}{Implicit and explicit definitions in \(L_{\omega_1\omega}\)}
Let \(L=\{0,S\}\) with \(0\) a constant symbol and \(S\) unary, write
\(S^{k}0\) for the term with \(k\) occurrences of \(S\), and let
\(\varphi(P)\) be the conjunction of
\(\forall x\,\bigvee_{k<\omega}x=S^{k}0\), of \(P(0)\), and of
\(\forall x\,(P(Sx)\leftrightarrow\neg P(x))\), with \(P\) unary.  In a
structure \(\mathfrak M\) for \(L\), a relation \(R\) with
\((\mathfrak M,R)\models\varphi(P)\) must contain \(0^{\mathfrak M}\) and
contain \((S^{k+1}0)^{\mathfrak M}\) exactly when it omits
\((S^{k}0)^{\mathfrak M}\); by induction on \(k\), whether
\((S^{k}0)^{\mathfrak M}\in R\) is fixed for every \(k\) (it holds iff
\(k\) is even), and by the first conjunct every element has this form, so
there is at most one such \(R\).  (When
\((S^{k}0)^{\mathfrak M}=(S^{m}0)^{\mathfrak M}\) with \(k-m\) odd there is
none.)  The formula \(\theta(x)=\bigvee_{k<\omega}x=S^{2k}0\) defines \(P\)
explicitly: in a model \((\mathfrak M,R)\), each element \(b\) equals
\((S^{k}0)^{\mathfrak M}\) for some \(k\), and for each such \(k\),
\(b\in R\) iff \(k\) is even; so \(b\in R\) iff \(b=(S^{2j}0)^{\mathfrak M}\)
for some \(j\).  Both \(\varphi\) and \(\theta\) use countably infinite
disjunctions.
\Cref{thm:p02-w02c-p02-003-beth} guarantees some explicit definition for
every implicit one; its proof produces one from an interpolant, which need
not be this \(\theta\).
\end{hbexample}

By part~1 of \cref{thm:p02-w02c-p02-003-wso-translation}, applied to
\(L\cup\{P\}\), a weak second-order sentence \(\varphi(P)\) that defines
\(P\) implicitly over the structures \(\mathrm{HF}(\mathfrak M)\) has a
translation \(\varphi^{\natural}\) in \(L_{\omega_1\omega}\) that defines
\(P\) implicitly, so \(P\) has an explicit definition in
\(L_{\omega_1\omega}\).  The failure of Beth's theorem that Barwise reports
for weak second-order logic is therefore the absence of an explicit
definition inside \(L^{*}\), not inside \(L_{\omega_1\omega}\).
\Cref{thm:p02-w02c-p02-003-beth} concerns a single sentence or a countable
set of sentences, as Makkai's statement concerns single sentences
\cite[ch.~A.7, \S1, p.~236]{Makkai1977}; an uncountable set of sentences of
\(L_{\omega_1\omega}\) is not a sentence of \(L_{\omega_1\omega}\), and the
proof above does not apply to it.

\emph{The sufficiency half of Barwise's remark on an infinitary rule.}
Barwise says that a Completeness Theorem for countable theories of
\(L_{\omega_1\omega}\) requires adding an infinitary rule of proof
\cite[ch.~A.1, \S5, p.~44]{Barwise1977FirstOrder}.
\Cref{thm:p02-w02c-p02-003-infinitary-rule} shows that some such rule is
needed.  Barwise does not say which rule, and he gives no calculus.
Makkai refers to a complete Gentzen-type infinitary system and attributes
the completeness theorem to Karp and Lopez-Escobar
\cite[ch.~A.7, \S1, p.~236]{Makkai1977}; neither source is held, and that
system is not reproduced here (its statement in the handbook is the cited
\cref{thm:p02-karp-completeness}).  The calculus below is the handbook's
own choice: Barwise's Hilbert-style system H
(\cref{def:p02-w02-p02-005-bw-system-h}) carried over to
\(L_{\omega_1\omega}\), working with sentences only, with new constant
symbols in the role of free variables, and with one rule that has
countably many premises.  The completeness theorem for it,
\cref{thm:p02-w02c-p02-003-infinitary-completeness}, is proved from a
model existence theorem for languages with function symbols,
\cref{thm:p02-w02c-p02-003-term-model-existence}, which generalizes
\cref{thm:p02-w02c-p02-003-model-existence}.

Throughout, formulas of \(L_{\omega_1\omega}\) are generated from atomic
formulas by \(\neg\), by \(\bigwedge\) and \(\bigvee\) of countable sets of
formulas, and by \(\exists x\), \(\forall x\); the finitary \(\land\),
\(\lor\) are \(\bigwedge\), \(\bigvee\) of two-element sets, \(\chi\to\chi'\)
is \(\neg\chi\lor\chi'\), \(\chi\leftrightarrow\chi'\) is
\((\chi\to\chi')\land(\chi'\to\chi)\), and \(\to\) associates to the right.
A sentence has no free variable.  A \emph{closed term} is a term without
variables.  For a formula \(\chi\) with at most \(x\) free and a closed term
\(t\), \(\chi(t/x)\) replaces every free occurrence of \(x\) by \(t\); as
\(t\) has no variable, no occurrence of a variable of \(t\) is captured.

\begin{hbdefinition}{def:p02-w02c-p02-003-hinf}{The calculus \(\mathbf H_{\infty}\) for \(L_{\omega_1\omega}\)}
Let \(L\) be a language in Barwise's sense, and let \(\mathrm{Par}\) be a
countably infinite set of constant symbols not in \(L\), the
\emph{parameters}.  Every line of \(\mathbf H_{\infty}\) is a sentence of
\(L_{\omega_1\omega}\) for \(L\cup\mathrm{Par}\).  The \emph{prime}
sentences are the atomic sentences and the sentences \(\exists x\,\chi\)
and \(\forall x\,\chi\).  A \emph{truth assignment} gives each prime
sentence a value \(\mathbf t\) or \(\mathbf f\); it extends to all
sentences by the clauses: \(\neg\chi\) has value \(\mathbf t\) iff
\(\chi\) has value \(\mathbf f\); \(\bigwedge\Phi\) has value \(\mathbf t\)
iff every member of \(\Phi\) has; \(\bigvee\Phi\) has value \(\mathbf t\)
iff some member of \(\Phi\) has.  A sentence is a \emph{tautology} if it
has value \(\mathbf t\) under every truth assignment.  The axioms of
\(\mathbf H_{\infty}\) are the following sentences, where \(t,t'\) are
closed terms of \(L\cup\mathrm{Par}\):
\begin{enumerate}
\item[(A1)] every tautology;
\item[(A2)] \(t=t\), and \(t=t'\to(\alpha(t/x)\to\alpha(t'/x))\) for every
      atomic formula \(\alpha\) with at most \(x\) free;
\item[(A3)] \(\forall x\,\chi\to\chi(t/x)\) and
      \(\chi(t/x)\to\exists x\,\chi\) for every formula \(\chi\) with at
      most \(x\) free.
\end{enumerate}
The rules are, for sentences \(\theta,\theta'\):
\begin{enumerate}
\item[(R1)] from \(\theta\) and \(\theta\to\theta'\) infer \(\theta'\);
\item[(R2)] from \(\theta\to\chi(p/x)\) infer \(\theta\to\forall x\,\chi\),
      and from \(\chi(p/x)\to\theta\) infer \(\exists x\,\chi\to\theta\),
      provided the parameter \(p\) occurs neither in \(\theta\) nor in
      \(\chi\);
\item[(R3)] from \(\theta\to\chi\) for every \(\chi\in\Phi\) infer
      \(\theta\to\bigwedge\Phi\), for every countable set \(\Phi\) of
      sentences.
\end{enumerate}
Let \(T\) be a set of sentences of \(L_{\omega_1\omega}\) for \(L\), so that
no parameter occurs in \(T\).  The set of \emph{theorems of \(T\)} is the
least set of sentences of \(L_{\omega_1\omega}\) for \(L\cup\mathrm{Par}\)
that contains every axiom and every member of \(T\) and is closed under
(R1)--(R3); \(T\vdash_{\infty}\theta\) says that \(\theta\) is a theorem of
\(T\).  (R3) is the infinitary rule: when \(\Phi\) is infinite it has
infinitely many premises.  (R2) mirrors the generalization rules of H,
with a parameter in place of a variable; (A3) is the quantifier axiom of H
restricted to closed terms, and (A2) plays the part of the equality axioms
of H in a substitution form.
\end{hbdefinition}

Since the theorems of \(T\) form the least set with these closure
properties, a property holds of all of them as soon as the set of
sentences with the property contains the axioms and \(T\) and is closed
under (R1)--(R3); proofs about \(\vdash_{\infty}\) below have this form.
Working with sentences and parameters avoids a difficulty that free
variables would raise: a single sentence of \(L_{\omega_1\omega}\) can
contain every variable of a countable stock, and then no variable is left
over to play the part of a new constant.

\begin{hblemma}{thm:p02-w02c-p02-003-hinf-soundness}{proved=annex:p02-w02c-p02-003-hinf-sound}{Soundness of \(\mathbf H_{\infty}\), and countably many premises}
Let \(T\) be a set of sentences of \(L_{\omega_1\omega}\) for \(L\).
\begin{enumerate}
\item If \(T\vdash_{\infty}\theta\), then every structure for
      \(L\cup\mathrm{Par}\) that satisfies every member of \(T\) satisfies
      \(\theta\).
\item Assume the axiom of choice.  If \(T\vdash_{\infty}\theta\), then
      \(T_0\vdash_{\infty}\theta\) for some countable \(T_0\subseteq T\).
\end{enumerate}
\end{hblemma}

\begin{hblemma}{thm:p02-w02c-p02-003-hinf-nnf}{proved=annex:p02-w02c-p02-003-hinf-nnf}{Derived rules and negation normal form in \(\mathbf H_{\infty}\)}
Let \(T\) be a set of sentences of \(L_{\omega_1\omega}\) for \(L\).
\begin{enumerate}
\item If \(T\vdash_{\infty}\theta_i\) for \(i=1,\ldots,n\) and
      \(\theta_1\to\cdots\to\theta_n\to\theta\) is a tautology, then
      \(T\vdash_{\infty}\theta\).
\item If \(\Psi\) is a countable set of sentences and
      \(T\vdash_{\infty}\chi\to\theta\) for every \(\chi\in\Psi\), then
      \(T\vdash_{\infty}\bigvee\Psi\to\theta\).  If \(T\) is countable,
      \(T\vdash_{\infty}\bigwedge T\).
\item Let \(\chi\) be a formula of \(L_{\omega_1\omega}\) for
      \(L\cup\mathrm{Par}\) with finitely many free variables in which only
      finitely many parameters occur, and let \(\sigma\) map the free
      variables of \(\chi\) to closed terms.  Write \(\chi\sigma\) for the
      sentence obtained by replacing each free occurrence of each such
      variable \(x\) by \(\sigma(x)\).  Then
      \(T\vdash_{\infty}\chi^{+}\sigma\leftrightarrow\chi\sigma\) and
      \(T\vdash_{\infty}\chi^{\ominus}\sigma\leftrightarrow\neg\chi\sigma\),
      with \(\chi^{+}\), \(\chi^{\ominus}\) as in
      \cref{def:p02-w02c-p02-003-hintikka}.
\end{enumerate}
\end{hblemma}

\begin{hbdefinition}{def:p02-w02c-p02-003-term-consistency}{Consistency property with closed terms}
Let \(L\) be a language and \(\mathrm{Par}\) a countably infinite set of
constant symbols not in \(L\).  A \emph{consistency property with closed
terms} for \(L\cup\mathrm{Par}\) is a set \(\mathcal S\) of finite sets of
n.n.f. sentences of \(L_{\omega_1\omega}\) for \(L\cup\mathrm{Par}\) such
that for every \(s\in\mathcal S\) and all closed terms \(t,t'\) of
\(L\cup\mathrm{Par}\):
\begin{enumerate}
\item[(T0)] for no atomic sentence \(\theta\) are both \(\theta\) and
      \(\neg\theta\) in \(s\);
\item[(T1)] \(s\cup\{t=t\}\in\mathcal S\);
\item[(T2)] if \(\alpha\) is an atomic formula with at most \(x\) free,
      \(\alpha(t/x)\in s\) and \((t=t')\in s\), then
      \(s\cup\{\alpha(t'/x)\}\in\mathcal S\);
\item[(T3)] if \(\bigwedge\Phi\in s\) and \(\varphi\in\Phi\), then
      \(s\cup\{\varphi\}\in\mathcal S\);
\item[(T4)] if \(\forall x\,\varphi\in s\), then
      \(s\cup\{\varphi(t/x)\}\in\mathcal S\);
\item[(T5)] if \(\bigvee\Phi\in s\), then \(s\cup\{\varphi\}\in\mathcal S\)
      for some \(\varphi\in\Phi\);
\item[(T6)] if \(\exists x\,\varphi\in s\), then
      \(s\cup\{\varphi(u/x)\}\in\mathcal S\) for some closed term \(u\).
\end{enumerate}
When \(L\) has no function symbols and its constant symbols are put into
\(\mathrm{Par}\), the closed terms are the constants, and (T0)--(T6) are
(C0)--(C6) of \cref{def:p02-w02c-p02-003-consistency} with (C6) weakened
to a witness among the constants.
\end{hbdefinition}

\begin{hbtheorem}{thm:p02-w02c-p02-003-term-model-existence}{proved=annex:p02-w02c-p02-003-term-mex}{Model existence with function symbols}
Assume the axiom of choice.  Let \(L\) be a countable language (relation,
function and constant symbols), \(\mathrm{Par}\) a countably infinite set of
constant symbols not in \(L\), and \(\mathcal S\) a consistency property
with closed terms for \(L\cup\mathrm{Par}\).  Then every
\(s\in\mathcal S\) has a model: a structure for \(L\cup\mathrm{Par}\) that
satisfies every member of \(s\), and in which every element is the value
of a closed term.
\end{hbtheorem}

\begin{hbtheorem}{thm:p02-w02c-p02-003-infinitary-completeness}{proved=annex:p02-w02c-p02-003-hinf-complete}{Completeness of \(\mathbf H_{\infty}\) for countable theories}
Assume the axiom of choice.  Let \(L\) be a language in Barwise's sense,
of any cardinality, \(T\) a countable set of sentences of
\(L_{\omega_1\omega}\) for \(L\), and \(\varphi\) a sentence of
\(L_{\omega_1\omega}\) for \(L\).  Then \(T\vdash_{\infty}\varphi\) if and
only if every structure for \(L\) that is a model of \(T\) satisfies
\(\varphi\).  In particular a sentence of \(L_{\omega_1\omega}\) is valid
iff it is a theorem of the empty set, and a countable set of sentences has
a model iff it does not prove \(\neg\bigwedge\varnothing\).
\end{hbtheorem}

Together with \cref{thm:p02-w02c-p02-003-infinitary-rule}, this gives
both halves of Barwise's remark for this calculus: completeness for
countable theories holds once the rule (R3) with countably many premises
is present, and fails for every calculus whose derivations use finitely
many premises.  The theorem concerns \(\mathbf H_{\infty}\); it is not a
proof of \cref{thm:p02-karp-completeness}, which is stated for Karp's
calculus.

\begin{hbexample}{def:p02-w02c-p02-003-hinf-example}{A derivation that uses the infinitary rule}
Let \(L\) contain a unary relation symbol \(R\) and constant symbols
\(c_0,c_1,\ldots\), and let
\(T=\{\bigvee_{n<\omega}R(c_n)\}\cup\{\neg R(c_n):n<\omega\}\).  Every
finite subset of \(T\) has a model (make \(R\) true of one constant not
mentioned among the negations), and \(T\) has none.  In
\(\mathbf H_{\infty}\), with \(\top=\bigwedge\varnothing\) (a tautology):
\(\neg R(c_n)\to(\top\to\neg R(c_n))\) is a tautology, so
\(T\vdash_{\infty}\top\to\neg R(c_n)\) for each \(n\) by (R1); (R3) gives
\(T\vdash_{\infty}\top\to\bigwedge_{n}\neg R(c_n)\), and (R1) with the axiom
\(\top\) gives \(T\vdash_{\infty}\bigwedge_n\neg R(c_n)\).  The sentence
\(\bigwedge_n\neg R(c_n)\to\bigvee_n R(c_n)\to\theta\) is a tautology for
every sentence \(\theta\) (no truth assignment makes both antecedents true),
so \(T\vdash_{\infty}\theta\) for every sentence \(\theta\), as the theorem
requires.  The one application of (R3) has infinitely many premises; by
\cref{thm:p02-w02c-p02-003-infinitary-rule} no calculus whose derivations
use finitely many members of \(T\) derives \(\neg\bigwedge\varnothing\)
from \(T\).
\end{hbexample}

\begin{hbtheorem}{thm:p02-w02c-p02-003-hinf-uncountable}{proved=annex:p02-w02c-p02-003-hinf-uncount}{Countability cannot be dropped}
Assume the axiom of choice.  Let \(L\) consist of constant symbols
\(c_\alpha\) for \(\alpha<\omega_1\) and \(d_n\) for \(n<\omega\), and let
\(T\) consist of \(\forall x\,\bigvee_{n<\omega}x=d_n\) and the sentences
\(\neg\,c_\alpha=c_\beta\) for \(\alpha<\beta<\omega_1\).  Then \(T\) has no
model, but \(T\nvdash_{\infty}\neg\bigwedge\varnothing\).
\end{hbtheorem}

The proof of \cref{thm:p02-w02c-p02-003-infinitary-completeness} uses
the countability of \(T\) in one place: it forms the sentence
\(\bigwedge T\).  \Cref{thm:p02-w02c-p02-003-hinf-uncountable} shows that
this is not an artefact of the proof: by part~2 of
\cref{thm:p02-w02c-p02-003-hinf-soundness}, every theorem of a theory is a
theorem of a countable part of it, so \(\mathbf H_{\infty}\) cannot detect
an inconsistency that only an uncountable set of premises produces.

The proof that Beth's theorem fails for weak second-order logic has
three steps: an \(L^{*}\)-sentence that pins down the standard model of
arithmetic, an implicit definition of arithmetical truth by Tarski's
clauses over that model, and a coding of \(\mathrm{HF}(\mathcal N)\) in
\(\mathcal N\) that turns an explicit weak second-order definition into
an arithmetical one.  The language is \(\Lang_A=\{0,S,+,\times\}\) and \(\mathcal N\)
is its standard structure (\cref{def:language-arithmetic}).

\begin{hbtheorem}{thm:p02-w02c-p02-003-wso-n}{proved=annex:p02-w02c-p02-003-wso-n}{A weak second-order sentence characterizing \(\mathcal N\)}
Let \(\alpha_{\mathcal N}\) be the \(L^{*}\)-sentence, for \(L=\Lang_A\),
that is the conjunction of the universal closures of the seven axioms of
\(\mathsf Q\) (\cref{sec:arithmetic-coding-01}) and of
\(\forall x\,\exists n\,\exists z\,(z=0\land\mathrm{Rec}_{\gamma_S}(n,z,x))\),
where the value sort is the individuals and \(\gamma_S(i,v,w)\) is
\(w=Sv\) (\cref{thm:p02-w02c-p02-003-hf-recursion}).  For every structure
\(\mathfrak M\) for \(\Lang_A\),
\(\mathrm{HF}(\mathfrak M)\models\alpha_{\mathcal N}\) iff
\(\mathfrak M\cong\mathcal N\).
\end{hbtheorem}

The remaining steps are an arithmetical coding of
\(\mathrm{HF}(\mathcal N)\) and an implicit truth definition.  Both use
primitive recursive functions (\cref{sec:turing-machines-02}) through their
graph formulas.  For a primitive recursive \(f\colon\N^{k}\to\N\) fix a
description \(d\) with \(f=f_d\) (\cref{thm:p02-w02b-p02-034-pr-graph}),
and for \(\Lang_A\)-terms \(r_1,\ldots,r_k,r\) write
\([f(r_1,\ldots,r_k)=r]\) for the formula obtained from
\(\Gamma_d(\bar x,y)\) of \cref{def:p05-pa-pr-graph} by renaming its free
variables to distinct variables occurring in none of the terms and then
substituting \(r_1,\ldots,r_k,r\) for them one at a time.  For a formula
\(\psi(y)\), \(\psi(f(r_1,\ldots,r_k))\) abbreviates
\(\exists y\,([f(r_1,\ldots,r_k)=y]\land\psi(y))\) with \(y\) new.  The
individual variables of \(L^{*}\) are \(x_0,x_1,\ldots\) and its set
variables are \(a_0,a_1,\ldots\).  For natural numbers \(x\) and \(y\geq1\),
\(\lfloor x/y\rfloor\) is the greatest \(q\) with \(qy\leq x\), and
\(\operatorname{pty}(x)=x\mathbin{\dot{-}}2\lfloor x/2\rfloor\) is the remainder of
\(x\) on division by \(2\).

\begin{hbdefinition}{def:p02-w02c-p02-003-hf-code}{Arithmetical codes of \(\mathrm{HF}(\mathcal N)\)}
The atoms of \(\mathrm{HF}(\mathcal N)\) are the natural numbers
(\cref{def:p02-w02-p02-005-bw-weak-so}).  Define
\(\operatorname{hc}_n\colon\N\cup\mathrm{HF}_n(\N)\to\N\) by
\(\operatorname{hc}_n(m)=2m\) for \(m\in\N\) and
\[
 \operatorname{hc}_{n+1}(a)=2\sum_{u\in a}2^{\operatorname{hc}_n(u)}+1
 \qquad(a\in\mathrm{HF}_{n+1}(\N)),
\]
and let \(\operatorname{hc}=\bigcup_n\operatorname{hc}_n\) (a function on
\(\N\cup\mathrm{HF}(\N)\) by part~1 of
\cref{thm:p02-w02c-p02-003-hf-arith}).  Let
\(\chi_{\in}(u,v)=\operatorname{pty}(v)\cdot
\operatorname{pty}(\lfloor\lfloor v/2\rfloor/2^{u}\rfloor)\), and let
\(\mathrm{Odd}(v)\) be the \(\Lang_A\)-formula
\(\exists w\,(v=S((SS0)\times w))\).  For an \(\Lang_A\)-term \(t\) of
\(L^{*}\), \(t^{\iota}\) is \(t\) with each \(x_i\) replaced by
\(x_{2i}\); for a set variable, \(a_i^{\iota}=x_{2i+1}\).  The
\emph{arithmetical translation} \(\theta^{\mathrm{ar}}\) of an
\(L^{*}\)-formula \(\theta\) for \(L=\Lang_A\) is defined by recursion on
\(\theta\):
\begin{align*}
 (r=t)^{\mathrm{ar}}&=(r^{\iota}=t^{\iota}),\\
 (t\in a_i)^{\mathrm{ar}}&=[\chi_{\in}((SS0)\times t^{\iota},x_{2i+1})=S0],\\
 (a_i\in a_j)^{\mathrm{ar}}&=[\chi_{\in}(x_{2i+1},x_{2j+1})=S0],\\
 (a_i=a_j)^{\mathrm{ar}}&=(x_{2i+1}=x_{2j+1}),\\
 (\exists x_i\,\theta)^{\mathrm{ar}}&=\exists x_{2i}\,\theta^{\mathrm{ar}},\\
 (\forall x_i\,\theta)^{\mathrm{ar}}&=\forall x_{2i}\,\theta^{\mathrm{ar}},\\
 (\exists a_i\,\theta)^{\mathrm{ar}}&=\exists x_{2i+1}\,(\mathrm{Odd}(x_{2i+1})\land\theta^{\mathrm{ar}}),\\
 (\forall a_i\,\theta)^{\mathrm{ar}}&=\forall x_{2i+1}\,(\mathrm{Odd}(x_{2i+1})\to\theta^{\mathrm{ar}}),
\end{align*}
and \(\theta^{\mathrm{ar}}\) commutes with every connective.  For an
assignment \(s\) in \(\mathrm{HF}(\mathcal N)\) (individual variables to
\(\N\), set variables to \(\mathrm{HF}(\N)\)), \(s^{\mathrm{ar}}\) is the
assignment in \(\mathcal N\) with \(s^{\mathrm{ar}}(x_{2i})=s(x_i)\) and
\(s^{\mathrm{ar}}(x_{2i+1})=\operatorname{hc}(s(a_i))\).
\end{hbdefinition}

\begin{hbtheorem}{thm:p02-w02c-p02-003-hf-arith}{proved=annex:p02-w02c-p02-003-hf-arith}{\(\mathrm{HF}(\mathcal N)\) is coded in \(\mathcal N\)}
\begin{enumerate}
\item \(\operatorname{hc}\) is a well-defined bijection of
      \(\N\cup\mathrm{HF}(\N)\) onto \(\N\); it maps the atoms onto the even
      numbers and \(\mathrm{HF}(\N)\) onto the odd numbers, and
      \(\operatorname{hc}(a)=2\sum_{u\in a}2^{\operatorname{hc}(u)}+1\) for
      \(a\in\mathrm{HF}(\N)\).
\item \(\chi_{\in}\) is primitive recursive, and for
      \(z\in\N\cup\mathrm{HF}(\N)\) and \(a\in\mathrm{HF}(\N)\),
      \(z\in a\) iff \(\chi_{\in}(\operatorname{hc}(z),\operatorname{hc}(a))=1\).
\item For every \(L^{*}\)-formula \(\theta\) for \(\Lang_A\) and every
      assignment \(s\) in \(\mathrm{HF}(\mathcal N)\),
      \(\mathrm{HF}(\mathcal N)\models\theta[s]\) iff
      \(\mathcal N\models\theta^{\mathrm{ar}}[s^{\mathrm{ar}}]\).
\item If \(\theta\) is an \(L^{*}\)-formula for \(\Lang_A\) whose free
      variables are among the individual variables
      \(x_{i_1},\ldots,x_{i_k}\), then the relation
      \(\{(m_1,\ldots,m_k):\mathrm{HF}(\mathcal N)\models\theta[s],\\
      s(x_{i_j})=m_j\}\) on \(\N\) is defined in \(\mathcal N\) by the
      first-order \(\Lang_A\)-formula \(\theta^{\mathrm{ar}}\), with
      \(x_{2i_j}\) in place of \(x_{i_j}\).
\end{enumerate}
\end{hbtheorem}

The truth clauses use the tree coding of \cref{sec:arithmetic-coding-04},
with its pairing \(C\), its projections \(h,t\) and its list entries
\((c)_j\).  Put \(\operatorname{arg}_0(c)=(t(c))_0\) and
\(\operatorname{arg}_1(c)=(t(c))_1\), so that \(h(c)\) is the outer tag of a
code \(c=N_k(c_0,c_1)\) and \(c_0,c_1\) are its arguments.  Let
\(\operatorname{Snt}(c)\) be \(1\) if \(c\) is the code of a first-order
\(\Lang_A\)-sentence and \(0\) otherwise; let
\(\operatorname{sb}(c,i,n)\) be the numeral substitution function of that
section, with
\(\operatorname{sb}(\ulcorner\varphi\urcorner,i,n)=\ulcorner\varphi[\overline n/x_i]\urcorner\);
and let \(\operatorname{tv}\) (term value) be defined by course of values:
\(\operatorname{tv}(c)\) is \(0\) if \(h(c)=1\),
\(\operatorname{tv}(\operatorname{arg}_0(c))+1\) if \(h(c)=2\),
\(\operatorname{tv}(\operatorname{arg}_0(c))+\operatorname{tv}(\operatorname{arg}_1(c))\)
if \(h(c)=3\),
\(\operatorname{tv}(\operatorname{arg}_0(c))\cdot\operatorname{tv}(\operatorname{arg}_1(c))\)
if \(h(c)=4\), and \(0\) otherwise.

\begin{hbdefinition}{def:p02-w02c-p02-003-truth-clauses}{Tarski's truth clauses for the codes of arithmetic sentences}
Let \(T\) be a new unary relation symbol.  \(\mathrm{TC}(T)\) is the
first-order sentence of \(\Lang_A\cup\{T\}\) that is the conjunction of
\(\forall c\,(T(c)\to[\operatorname{Snt}(c)=\overline1])\) and, for each
tag \(k\) from \(5\) to \(12\), of
\(\forall c\,(([\operatorname{Snt}(c)=\overline1]\land[h(c)=\overline k])\to\beta_k(c))\),
where
\begin{align*}
 \beta_5(c)&:\ T(c)\leftrightarrow
   \operatorname{tv}(\operatorname{arg}_0(c))=\operatorname{tv}(\operatorname{arg}_1(c)),\\
 \beta_6(c)&:\ \neg T(c),\\
 \beta_7(c)&:\ T(c)\leftrightarrow\neg T(\operatorname{arg}_0(c)),\\
 \beta_8(c)&:\ T(c)\leftrightarrow(T(\operatorname{arg}_0(c))\land T(\operatorname{arg}_1(c))),\\
 \beta_9(c)&:\ T(c)\leftrightarrow(T(\operatorname{arg}_0(c))\lor T(\operatorname{arg}_1(c))),\\
 \beta_{10}(c)&:\ T(c)\leftrightarrow(T(\operatorname{arg}_0(c))\to T(\operatorname{arg}_1(c))),\\
 \beta_{11}(c)&:\ T(c)\leftrightarrow\forall n\,T(\operatorname{sb}(\operatorname{arg}_1(c),\operatorname{arg}_0(c),n)),\\
 \beta_{12}(c)&:\ T(c)\leftrightarrow\exists n\,T(\operatorname{sb}(\operatorname{arg}_1(c),\operatorname{arg}_0(c),n)),
\end{align*}
with the function terms read by the abbreviation above.  The tags
\(5,\ldots,12\) are those of \(r=t\), \(\bottom\), \(\neg\), \(\land\),
\(\lor\), \(\to\), \(\forall\), \(\exists\) in that coding.
\(\alpha_{\mathrm{TA}}\) is the \(L^{*}\)-sentence
\(\alpha_{\mathcal N}\land\mathrm{TC}(T)\) for \(\Lang_A\cup\{T\}\).
\end{hbdefinition}

\begin{hblemma}{thm:p02-w02c-p02-003-truth-clauses}{proved=annex:p02-w02c-p02-003-truth-clauses}{The truth clauses have true arithmetic as their only solution over \(\mathcal N\)}
\begin{enumerate}
\item \(h\), \(\operatorname{arg}_0\), \(\operatorname{arg}_1\),
      \(\operatorname{Snt}\), \(\operatorname{sb}\) and \(\operatorname{tv}\)
      are primitive recursive, and \(\operatorname{tv}(\ulcorner r\urcorner)\)
      is the value \(r^{\mathcal N}\) of every closed \(\Lang_A\)-term \(r\).
\item For every \(X\subseteq\N\), \((\mathcal N,X)\models\mathrm{TC}(T)\)
      iff \(X=\mathrm{TA}\).
\end{enumerate}
\end{hblemma}

\begin{hbtheorem}{thm:p02-w02c-p02-003-wso-beth-failure}{proved=annex:p02-w02c-p02-003-wso-beth-failure}{Beth's theorem fails for weak second-order logic (Barwise, p.~44)}
\begin{enumerate}
\item \(\alpha_{\mathrm{TA}}\) defines \(T\) implicitly: for every structure
      \(\mathfrak M\) for \(\Lang_A\) there is at most one \(X\subseteq M\)
      with \(\mathrm{HF}((\mathfrak M,X))\models\alpha_{\mathrm{TA}}\); and
      \(\mathrm{HF}((\mathcal N,\mathrm{TA}))\models\alpha_{\mathrm{TA}}\).
\item \(T\) is not explicitly definable in weak second-order logic: there is
      no \(L^{*}\)-formula \(\theta(x)\) for \(\Lang_A\), whose only free
      variable is the individual variable \(x\), such that every model
      \(\mathrm{HF}((\mathfrak M,X))\) of \(\alpha_{\mathrm{TA}}\) satisfies
      \(\forall x\,(T(x)\leftrightarrow\theta(x))\).
\end{enumerate}
\end{hbtheorem}

Barwise names no relation; \(T\) interpreted as \(\mathrm{TA}\) is one
witness.  By the paragraph after \cref{def:p02-w02c-p02-003-beth-example},
the translation \(\alpha_{\mathrm{TA}}^{\natural}\) defines \(T\)
implicitly in \(L_{\omega_1\omega}\), so \(T\) has an explicit definition
there; one is \(\bigvee\{x=\overline n:n\in\mathrm{TA}\}\), a countable
disjunction.  The proof of part~2 uses
\cref{thm:p02-w02c-p02-003-hf-arith} to turn an explicit weak second-order
definition into an arithmetical truth definition, which
\cref{thm:tarski-undefinability} excludes.

Barwise's claim that \(L(Q_0)\) is essentially equivalent to
\(\omega\)-logic and weak second-order logic
\cite[ch.~A.1, \S5, p.~44]{Barwise1977FirstOrder} comes without a precise
sense of ``essentially equivalent''.  The comparison below is made for
classes of structures that are reducts of relativized models, so that each
logic may use new symbols and new elements.

\begin{hbconvention}{def:p02-w02c-p02-003-rpc}{Relativized projective classes and essential equivalence}
Let \(L\) be a language, \(U\) a unary relation symbol not in \(L\), and
\(L'\supseteq L\cup\{U\}\).  For a structure \(\mathfrak A\) for \(L'\) in
which \(U^{\mathfrak A}\) is nonempty, contains \(c^{\mathfrak A}\) for each
constant symbol \(c\in L\) and is closed under \(f^{\mathfrak A}\) for each
function symbol \(f\in L\), \(\mathfrak A^{(U)}_{L}\) is the structure for
\(L\) with domain \(U^{\mathfrak A}\) and the restrictions to it of the
interpretations of the symbols of \(L\).
\begin{enumerate}
\item The three logics have the following sentences and models for a
      language \(L'\).  \(L(Q_0)\): its sentences for \(L'\), with
      \(\mathfrak A\models\varphi\) as in
      \cref{def:p02-w02-p02-005-bw-infinitary}.  Weak second-order logic:
      the sentences \(\varphi\) of \(L'^{*}\), with \(\mathfrak A\) a model
      of \(\varphi\) iff \(\mathrm{HF}(\mathfrak A)\models\varphi\).
      \(\omega\)-logic: the first-order sentences of an
      \(L'\supseteq\{\mathrm{Nn},\hat0,\hat S,\hat+,\hat\times\}\), with
      \(\mathrm{Nn}\) unary, \(\hat0\) a constant symbol, \(\hat S\) unary and
      \(\hat+,\hat\times\) binary function symbols, whose models are the
      \emph{\(\omega\)-models} satisfying it: structures \(\mathfrak A\)
      for \(L'\) in which \(\mathrm{Nn}^{\mathfrak A}\) contains
      \(\hat0^{\mathfrak A}\), is closed under \(\hat S^{\mathfrak A}\),
      \(\hat+^{\mathfrak A}\), \(\hat\times^{\mathfrak A}\), and with the
      restricted operations is isomorphic to \(\mathcal N\).
\item A class \(K\) of structures for \(L\) is \emph{RPC} in one of these
      logics if for some \(L'\supseteq L\cup\{U\}\) there is a sentence
      \(\varphi\) of that logic for \(L'\) such that every model
      \(\mathfrak A\) of \(\varphi\) has \(U^{\mathfrak A}\) nonempty and
      closed as above, and \(K\) is the class of structures isomorphic to
      some \(\mathfrak A^{(U)}_{L}\) with \(\mathfrak A\) a model of
      \(\varphi\).
\item Two of the logics are \emph{essentially equivalent} if, for every
      language \(L\), they have the same RPC classes of structures for
      \(L\).
\end{enumerate}
The \(\omega\)-logic of item~1 is the one-sorted form
\((M\cup N,M,N,\ldots)\) of Barwise's two-sorted structures
\((\mathfrak M,\mathbb N,\ldots)\) recorded in
\cref{def:p02-w02-p02-005-bw-many-sorted}, with the number sort the
predicate \(\mathrm{Nn}\); Barwise does not say which operations the fixed
structure \(\mathbb N\) carries, and item~1 takes those of \(\mathcal N\).
\end{hbconvention}

\begin{hbtheorem}{thm:p02-w02c-p02-003-q0-equivalence}{proved=annex:p02-w02c-p02-003-q0-equivalence}{\(L(Q_0)\), \(\omega\)-logic and weak second-order logic are essentially equivalent (Barwise, p.~44)}
Assume the axiom of choice, and let \(L\) be a language.
\begin{enumerate}
\item For every formula \(\varphi\) of \(L(Q_0)\) for \(L\) there is a
      formula \(\varphi^{w}\) of \(L^{*}\) with the same free variables such
      that \(\mathfrak M\models\varphi[s]\) iff
      \(\mathrm{HF}(\mathfrak M)\models\varphi^{w}[s]\) for every structure
      \(\mathfrak M\) for \(L\) and assignment \(s\) in \(M\): \(\varphi^{w}\)
      replaces each subformula \(Q_0x\,\psi\) by
      \(\neg\exists a\,\forall x\,(\psi^{w}\to x\in a)\), with \(a\) a set
      variable not in \(\psi^{w}\).  So every class RPC in \(L(Q_0)\) is RPC
      in weak second-order logic.
\item Every class RPC in weak second-order logic is RPC in \(L(Q_0)\).
\item Every class RPC in \(L(Q_0)\) is RPC in \(\omega\)-logic.
\item Every class RPC in \(\omega\)-logic is RPC in weak second-order
      logic.
\end{enumerate}
Hence the three logics are essentially equivalent in the sense of
\cref{def:p02-w02c-p02-003-rpc}.  The axiom of choice is used only in
part~3.
\end{hbtheorem}

The proof of part~2 describes \(\mathrm{HF}(\mathfrak M)\) inside a single
structure by first-order axioms together with two uses of \(Q_0\): every
set has finitely many members, and a rank function into a linear order in
which each element has finitely many predecessors makes membership
well founded.  Part~3 replaces each subformula \(Q_0x\,\psi\) by a new
relation symbol and adds functions that witness its truth (an injection of
the number sort) or its falsity (an injection into an initial segment of
the number sort).  Part~4 pins down the number sort in weak second-order
logic with \cref{thm:p02-w02c-p02-003-hf-recursion}.  Only part~1 is a
translation sentence by sentence over the same language; parts~2--4 use
new symbols and, in parts~2 and~3, new elements.

\subsection{Keisler's Completeness and Compactness Theorems for \(L(Q)\)}
\label{sec:p02-w02c-p02-003-keisler}

\emph{Purpose.}  Barwise reports that \(L(Q)\), with \(Q\) read as ``there
exist uncountably many'', has a Completeness Theorem and a Compactness
Theorem over a finite or countable language, analogous to those of
first-order logic, with the proof due to Keisler
\cite[ch.~A.1, \S5, pp.~44--45]{Barwise1977FirstOrder}.  This subsection
states Keisler's axioms and rules and proves both theorems, following
Keisler's paper \cite[\S\S1--2, pp.~4--27]{Keisler1970}.  \emph{Inputs}: a
first-order language \(L\) with identity and at most countably many
predicate, function and individual constant symbols (Keisler's standing
hypothesis, p.~4, where ``countable'' means of power at most \(\omega\)),
and a set \(\Sigma\) of sentences of \(L(Q)\).  \emph{Conclusions}:
\(\Sigma\) has a standard model iff it is consistent in \(L(Q)\); if every
finite subset of \(\Sigma\) has a standard model, so does \(\Sigma\).  The
proof has two stages: a Henkin construction gives a countable \emph{weak}
model, in which \(Q\) is interpreted by an arbitrary family of sets; an
elementary chain of length \(\omega_1\) of countable weak models then has a
union in which \(Q\) means ``uncountably many''.  Keisler writes the
identity symbol as \(\equiv\); here it is written \(=\), since \(\equiv\)
has other registered roles in the handbook.

\begin{hbdefinition}{def:p02-w02c-p02-003-kq-language}{\(L(Q)\), weak models and standard models (Keisler, \S1)}
The formulas of \(L(Q)\) form the least set containing the atomic formulas
of \(L\) and closed under: if \(\varphi,\psi\) are formulas and \(y\) is a
variable, then \(\varphi\land\psi\), \(\neg\varphi\), \((\exists y)\varphi\),
\((\forall y)\varphi\) and \((Qy)\varphi\) are formulas
\cite[\S1, p.~4]{Keisler1970}.  \(\varphi\lor\psi\), \(\varphi\to\psi\) and
\(\varphi\leftrightarrow\psi\) abbreviate their usual definitions from
\(\neg\) and \(\land\).  \((Qy)\) binds \(y\) as \((\exists y)\) does; free
and bound occurrences, sentences, \(\varphi(t/x)\) (replacement of the free
occurrences of \(x\) by the term \(t\)) and ``\(t\) is substitutable for
\(x\) in \(\varphi\)'' (no free occurrence of \(x\) lies in the scope of
\((\exists y)\), \((\forall y)\) or \((Qy)\) with \(y\) in \(t\)) are
defined with this clause added.  \(\varphi(v_1,\ldots,v_n)\) denotes a
formula whose free variables are among \(v_1,\ldots,v_n\).
\begin{enumerate}
\item A \emph{weak model} is a pair \((\mathfrak A,q)\) of a structure
      \(\mathfrak A\) for \(L\), with universe \(A\neq\varnothing\), and a set
      \(q\subseteq S(A)\) of subsets of \(A\).  Satisfaction
      \((\mathfrak A,q)\models\varphi[a_1,\ldots,a_n]\) is defined by
      induction on \(\varphi\) with the first-order clauses and the clause:
      \((\mathfrak A,q)\models(Qv_m)\psi[a_1,\ldots,a_n]\) iff the set of
      \(b\in A\) with \((\mathfrak A,q)\models\psi[a_1,\ldots,a_{m-1},b,
      a_{m+1},\ldots,a_n]\) is a member of \(q\)
      \cite[\S1, pp.~4--5]{Keisler1970}.  \((\mathfrak A,q)\models\varphi\)
      means that the universal closure of \(\varphi\) holds.
\item \(\mathfrak A\models\varphi[a_1,\ldots,a_n]\) means
      \((\mathfrak A,q)\models\varphi[a_1,\ldots,a_n]\) with \(q\) the set of
      all uncountable subsets of \(A\); \(\mathfrak A\) is a \emph{standard
      model} of a sentence \(\varphi\) iff \(\mathfrak A\models\varphi\)
      \cite[\S1, p.~5]{Keisler1970}.  This is the semantics of \(L(Q)\) in
      \cref{def:p02-w02-p02-005-bw-infinitary}.
\item For a set \(\Sigma(x)\) of formulas with at most \(x\) free,
      \((\mathfrak A,q)\) \emph{omits} \(\Sigma\) iff no \(a\in A\) satisfies
      every member of \(\Sigma\) in \((\mathfrak A,q)\)
      \cite[\S1, p.~5]{Keisler1970}.
\item \((\mathfrak B,r)\) is an \emph{elementary extension} of
      \((\mathfrak A,q)\), written \((\mathfrak A,q)\prec(\mathfrak B,r)\),
      iff \(A\subseteq B\) and for all formulas \(\varphi(x_1,\ldots,x_n)\)
      of \(L(Q)\) and \(a_1,\ldots,a_n\in A\),
      \((\mathfrak A,q)\models\varphi[a_1,\ldots,a_n]\) iff
      \((\mathfrak B,r)\models\varphi[a_1,\ldots,a_n]\).  A sequence
      \((\mathfrak A_\alpha,q_\alpha)\), \(\alpha<\gamma\), is an
      \emph{elementary chain} iff
      \((\mathfrak A_\alpha,q_\alpha)\prec(\mathfrak A_\beta,q_\beta)\) for
      \(\alpha<\beta<\gamma\); its \emph{union}
      \((\mathfrak A,q)=\bigcup_{\alpha<\gamma}(\mathfrak A_\alpha,q_\alpha)\)
      has \(\mathfrak A=\bigcup_{\alpha<\gamma}\mathfrak A_\alpha\) and
      \(q=\{S\subseteq A:\text{for some }\beta<\gamma,\ \beta\le\alpha<\gamma
      \text{ implies }S\cap A_\alpha\in q_\alpha\}\)
      \cite[\S2, p.~16]{Keisler1970}.
\item For a structure \(\mathfrak A\), \(\mathfrak A^{*}\) is its expansion
      by a new constant \(c_a\) for each \(a\in A\), interpreted by \(a\);
      \(L^{*}\) is the language of \(\mathfrak A^{*}\)
      \cite[\S2, p.~18]{Keisler1970}.
\end{enumerate}
\end{hbdefinition}

\begin{hbdefinition}{def:p02-w02c-p02-003-kq-calculus}{Keisler's axioms and rules for \(L(Q)\)}
The axioms of \(L(Q)\) are the formulas of the following schemes, where
\(\varphi,\psi\) are formulas of \(L(Q)\), which may have free variables,
and \(x,y,z\) are distinct variables \cite[\S1, p.~6]{Keisler1970}.
\begin{description}
\item[Axiom 0.] The axiom schemes of first-order logic with identity,
      instantiated by formulas of \(L(Q)\).  Keisler takes them from
      Church's textbook; the handbook fixes the following list:
      (0a) every formula obtained from a tautology of propositional logic in
      \(\neg,\land\) by substituting formulas of \(L(Q)\) for its sentence
      letters; (0b) \((\forall x)\varphi\to\varphi(t/x)\) and (0c)
      \(\varphi(t/x)\to(\exists x)\varphi\), with \(t\) substitutable for
      \(x\) in \(\varphi\); (0d) \((\forall x)(\psi\to\varphi)\to(\psi\to(\forall
      x)\varphi)\) and (0e) \((\forall x)(\varphi\to\psi)\to((\exists
      x)\varphi\to\psi)\), with \(x\) not free in \(\psi\); (0f) \(x=x\); (0g)
      \(x=y\to(\alpha\to\alpha')\), for \(\alpha\) atomic and \(\alpha'\)
      obtained from \(\alpha\) by replacing some occurrences of \(x\) by
      \(y\).
\item[Axiom 1.] \(\neg(Qx)(x=y\lor x=z)\).
\item[Axiom 2.] \((\forall x)(\varphi\to\psi)\to((Qx)\varphi\to(Qx)\psi)\).
\item[Axiom 3.] \((Qx)\varphi(x\ldots)\leftrightarrow(Qy)\varphi(y\ldots)\),
      where \(y\) does not occur in \(\varphi(x\ldots)\) and
      \(\varphi(y\ldots)\) comes from it by replacing each free occurrence of
      \(x\) by \(y\).
\item[Axiom 4.] \((Qy)(\exists x)\varphi\to(\exists x)(Qy)\varphi\lor(Qx)(\exists y)\varphi\).
\end{description}
The rules are modus ponens (from \(\varphi\) and \(\varphi\to\psi\) infer
\(\psi\)) and generalization (from \(\varphi\) infer \((\forall x)\varphi\))
\cite[\S1, p.~7]{Keisler1970}.  \(\vdash\varphi\) means that \(\varphi\) is
provable.  For a set \(\Sigma\) of \emph{sentences}, \(\Sigma\vdash\varphi\)
means that some finite sequence ending in \(\varphi\) consists of axioms,
members of \(\Sigma\) and results of the two rules applied to earlier
members; deductions are taken only from sets of sentences
\cite[\S1, p.~7]{Keisler1970}.  \(\Sigma\) is \emph{consistent} iff no
sentence \(\sigma\) has \(\Sigma\vdash\sigma\) and \(\Sigma\vdash\neg\sigma\).
\end{hbdefinition}

Keisler reads Axiom~1 as ``every set of power at most \(2\) is
countable'', Axiom~2 as ``every set with an uncountable subset is
uncountable'' and Axiom~4 as ``the union of countably many countable sets
is countable'' \cite[\S1, p.~6]{Keisler1970}.  If \(q\) is the set of all
infinite subsets of \(A\), every axiom holds in \((\mathfrak A,q)\)
\cite[\S1, p.~7]{Keisler1970}; so the axioms do not force \(Q\) to mean
``uncountably many'' in a weak model, and the second stage of the proof is
needed.

\begin{hblemma}{thm:p02-w02c-p02-003-kq-soundness}{proved=annex:p02-w02c-p02-003-kq-soundness}{Substitution and soundness (Keisler, 1.1--1.5)}
\begin{enumerate}
\item (Coincidence and substitution.)  Satisfaction of \(\varphi\) in a weak
      model depends only on the values of the free variables of
      \(\varphi\).  If \(t\) is substitutable for \(x\) in \(\varphi\), then
      \((\mathfrak A,q)\models\varphi(t/x)[s]\) iff
      \((\mathfrak A,q)\models\varphi[s(t^{\mathfrak A}[s]/x)]\).
\item Every instance of Axiom~0 and of Axiom~3 holds in every weak model.
\item If \((\mathfrak A,q)\) is a weak model of all the axioms of
      \(L(Q)\), \(\Sigma\vdash\varphi\) and \((\mathfrak A,q)\models\Sigma\),
      then \((\mathfrak A,q)\models\varphi\).
\item Every structure \(\mathfrak A\) for \(L\) is a standard model of
      every axiom.  Hence if \(\Sigma\) has a standard model, \(\Sigma\) is
      consistent, and if \(\mathfrak A\) is a standard model of \(\Sigma\)
      and \(\Sigma\vdash\varphi\), then \(\mathfrak A\models\varphi\).
\end{enumerate}
\end{hblemma}

Keisler remarks that the axiom of choice enters only in showing that
Axiom~4 holds in all standard models \cite[\S1, p.~7]{Keisler1970}; the
proof of part~4 uses it in the form that a countable union of countable
sets is countable.

\begin{hblemma}{thm:p02-w02c-p02-003-kq-proof-lemmas}{proved=annex:p02-w02c-p02-003-kq-proof-lemmas}{Proof-theoretic lemmas (Keisler, 1.6--1.8)}
Let \(\Sigma\) be a set of sentences of \(L(Q)\).
\begin{enumerate}
\item (Deduction theorem.)  If \(\Sigma\cup\{\sigma\}\vdash\varphi\) with
      \(\sigma\) a sentence, then \(\Sigma\vdash\sigma\to\varphi\).
\item (Derived rules.)  If \(\Sigma\vdash\psi\to\varphi\) and \(x\) is not
      free in \(\psi\), then \(\Sigma\vdash\psi\to(\forall x)\varphi\); if
      \(\Sigma\vdash\varphi\to\psi\) and \(x\) is not free in \(\psi\), then
      \(\Sigma\vdash(\exists x)\varphi\to\psi\).  If
      \(\Sigma\vdash\varphi\leftrightarrow\varphi'\), then
      \(\Sigma\vdash(Kx)\varphi\leftrightarrow(Kx)\varphi'\) for
      \(K\in\{\exists,\forall,Q\}\).
\item (Alphabetic variants.)  If \(\varphi'\) comes from \(\varphi\) by
      replacing a subformula \((Ky)\psi\) by \((Kz)\psi(z/y)\), with
      \(z\) not occurring in \(\psi\), then \(\vdash\varphi\leftrightarrow\varphi'\).
\item (New constants.)  Let \(L'\supseteq L\) add new individual
      constants.  If \(\Sigma\vdash\chi\) in \(L'(Q)\) and
      \(c_1,\ldots,c_k\) are new constants not occurring in \(\Sigma\), then
      for all distinct variables \(v_1,\ldots,v_k\) not occurring in \(\chi\),
      \(\Sigma\vdash\chi(v_1/c_1,\ldots,v_k/c_k)\), the result of replacing
      each occurrence of \(c_i\) by \(v_i\).  In particular \(\Sigma\) is
      consistent in \(L(Q)\) iff it is consistent in \(L'(Q)\).
\item (Maximal sets.)  If \(\Gamma\) is a maximal consistent set of
      sentences, then \(\Gamma\vdash\sigma\) iff \(\sigma\in\Gamma\),
      \(\neg\sigma\in\Gamma\) iff \(\sigma\notin\Gamma\), and
      \(\sigma\land\tau\in\Gamma\) iff \(\sigma\in\Gamma\) and \(\tau\in\Gamma\).
\end{enumerate}
\end{hblemma}

\begin{hbtheorem}{thm:p02-w02c-p02-003-kq-weak-completeness}{proved=annex:p02-w02c-p02-003-kq-weak-completeness}{Weak completeness and omitting types (Keisler, 2.2--2.4)}
Let \(L\) be finite or countable and \(\Sigma\) a set of sentences of
\(L(Q)\).
\begin{enumerate}
\item (Keisler's Lemma~2.4.)  For each \(n\) in a set \(N\subseteq\omega\)
      let \(\Sigma_n(x_n)\) be a set of formulas of \(L(Q)\) with only \(x_n\)
      free.  Suppose \(\Sigma\) is consistent and, for every \(n\in N\) and
      every formula \(\varphi(x_n)\) of \(L(Q)\), if \((\exists x_n)\varphi\)
      is consistent with \(\Sigma\), then some \(\sigma\in\Sigma_n\) makes
      \((\exists x_n)(\varphi\land\neg\sigma)\) consistent with \(\Sigma\).
      Then \(\Sigma\) has a countable weak model in which all the axioms of
      \(L(Q)\) hold and which omits each \(\Sigma_n\), \(n\in N\).
\item (Weak completeness, Keisler's Lemma~2.3.)  \(\Sigma\) is consistent
      iff it has a countable weak model in which all the axioms of
      \(L(Q)\) hold.
\item For every sentence \(\sigma\), \(\Sigma\vdash\sigma\) iff \(\sigma\)
      holds in every weak model of \(\Sigma\) in which all the axioms of
      \(L(Q)\) hold.
\end{enumerate}
\end{hbtheorem}

\begin{hblemma}{thm:p02-w02c-p02-003-kq-finite-unions}{proved=annex:p02-w02c-p02-003-kq-finite-unions}{Finite unions of small definable sets (semantic form of Keisler, 1.9)}
Let \((\mathfrak A,q)\) be a weak model in which all the axioms of \(L(Q)\)
hold, and call \(X\subseteq A\) \emph{definable} if
\(X=\{b:(\mathfrak A,q)\models\varphi[b,a_1,\ldots,a_n]\}\) for some
formula \(\varphi\) of \(L(Q)\) and \(a_1,\ldots,a_n\in A\).  For definable
\(X,Y\): if \(X\subseteq Y\) and \(X\in q\), then \(Y\in q\); no set with at
most two elements is in \(q\), and \(\varnothing\notin q\); if
\(X\cup Y\in q\), then \(X\in q\) or \(Y\in q\).  Consequently a definable
set included in a finite union of definable sets not in \(q\) is not in
\(q\).
\end{hblemma}

\begin{hblemma}{thm:p02-w02c-p02-003-kq-chains}{proved=annex:p02-w02c-p02-003-kq-chains}{Elementary chains and named elements (Keisler, 2.5--2.6)}
\begin{enumerate}
\item The union \((\mathfrak A,q)\) of an elementary chain
      \((\mathfrak A_\alpha,q_\alpha)\), \(\alpha<\gamma\), of weak models is
      a weak model and \((\mathfrak A_\alpha,q_\alpha)\prec(\mathfrak A,q)\)
      for every \(\alpha<\gamma\).
\item If all the axioms of \(L(Q)\) hold in \((\mathfrak A,q)\), all the
      axioms of \(L^{*}(Q)\) hold in \((\mathfrak A^{*},q)\).
\item If \((\mathfrak A,q)\prec(\mathfrak B,r)\), then
      \((\mathfrak A^{*},q)\prec(\mathfrak B',r)\), where \(\mathfrak B'\)
      expands \(\mathfrak B\) by interpreting each \(c_a\) as \(a\).
\end{enumerate}
\end{hblemma}

\begin{hblemma}{thm:p02-w02c-p02-003-kq-main-lemma}{proved=annex:p02-w02c-p02-003-kq-main-lemma}{Keisler's Main Lemma (2.7) and its iteration (2.8)}
Let \(L\) be finite or countable, let \((\mathfrak A,q)\) be a countable
weak model in which all the axioms of \(L(Q)\) hold, and let \(L^{*}\) be
the language of \(\mathfrak A^{*}\).
\begin{enumerate}
\item If \(\varphi(x)\) is a formula of \(L^{*}(Q)\) with
      \((\mathfrak A^{*},q)\models(Qx)\varphi(x)\), there is a countable
      elementary extension \((\mathfrak B,r)\) of \((\mathfrak A,q)\) such
      that (i) \((\mathfrak B^{*},r)\models\varphi[b]\) for some
      \(b\in B-A\), and (ii) for every formula \(\psi(y)\) of \(L^{*}(Q)\)
      with \((\mathfrak A^{*},q)\models\neg(Qy)\psi(y)\),
      \(\{a\in B:(\mathfrak B^{*},r)\models\psi[a]\}\subseteq A\).
      Here \(\mathfrak B^{*}\) interprets each \(c_a\) as \(a\).
\item \((\mathfrak A,q)\) has a countable elementary extension
      \((\mathfrak B,r)\) such that, for all formulas \(\varphi(x)\) of
      \(L^{*}(Q)\), \((\mathfrak A^{*},q)\models(Qx)\varphi(x)\) iff some
      \(b\in B-A\) has \((\mathfrak B^{*},r)\models\varphi[b]\).
\end{enumerate}
\end{hblemma}

\begin{hbtheorem}{thm:p02-w02c-p02-003-keisler-completeness}{proved=annex:p02-w02c-p02-003-keisler-completeness}{Keisler's Completeness Theorem for \(L(Q)\)}
Let \(L\) be a first-order language with identity and at most countably
many predicate, function and individual constant symbols, and let
\(\Sigma\) be a set of sentences of \(L(Q)\).  Then \(\Sigma\) has a
standard model iff \(\Sigma\) is consistent in \(L(Q)\) with the axioms and
rules of \cref{def:p02-w02c-p02-003-kq-calculus}
\cite[\S2, p.~13]{Keisler1970}.  If \(\Sigma\) is consistent, it has a
standard model of power at most \(\omega_1\); if moreover
\(\Sigma\cup\{(Qx)(x=x)\}\) is consistent, it has one of power exactly
\(\omega_1\) \cite[\S2, Corollary~2.9, p.~26]{Keisler1970}.  A sentence of
\(L(Q)\) is provable iff it holds in every standard model
\cite[\S2, Corollary~2.10, p.~26]{Keisler1970}.
\end{hbtheorem}

\begin{hbcorollary}{thm:p02-w02c-p02-003-keisler-compactness}{proved=annex:p02-w02c-p02-003-keisler-compactness}{Keisler's Compactness Theorem for \(L(Q)\)}
Let \(L\) be as in \cref{thm:p02-w02c-p02-003-keisler-completeness} and
\(\Sigma\) a set of sentences of \(L(Q)\).  If every finite subset of
\(\Sigma\) has a standard model, then \(\Sigma\) has a standard model
\cite[\S2, Corollary~2.11, p.~27]{Keisler1970}.
\end{hbcorollary}

Keisler records that compactness for \(L(Q)\) was first proved by Fuhrken,
and the recursive enumerability of the valid sentences by Vaught, by other
methods \cite[\S2, p.~27]{Keisler1970}.  The sentence \((Qx)(x=x)\) is not
an axiom: standard models may be countable, and in a countable one every
\((Qx)\)-formula is false \cite[\S1, p.~6]{Keisler1970}.

\begin{hbexample}{def:p02-w02c-p02-003-kq-example}{Failure for uncountably many constants}
Let \(L\) have a unary predicate \(P\) and constants \(c_\alpha\),
\(\alpha<\omega_1\), and let
\[
 \Delta=\{\neg(Qx)P(x)\}\cup\{P(c_\alpha):\alpha<\omega_1\}
        \cup\{\neg\,c_\alpha=c_\beta:\alpha<\beta<\omega_1\}
\]
\cite[\S2, p.~27]{Keisler1970}.  Every finite subset of \(\Delta\) has a
standard model: it names finitely many constants \(c_{\alpha_1},\ldots,
c_{\alpha_k}\); take universe \(\omega\), interpret these by \(1,\ldots,k\),
every other constant by \(0\), and \(P\) as \(\{1,\ldots,k\}\), a finite
and hence countable set.  \(\Delta\) itself has no standard model, since in
any model of the last two groups \(P\) contains the \(\omega_1\) distinct
values of the \(c_\alpha\).  So the Compactness Theorem fails for this
language, and so does the Completeness Theorem: by part~4 of
\cref{thm:p02-w02c-p02-003-kq-soundness} every finite subset of \(\Delta\)
is consistent, and a deduction uses finitely many members, so \(\Delta\) is
consistent without a standard model.  The countability hypothesis is used
in the proof at the omitting-types step and in the Main Lemma, where the
formulas of \(L^{*}(Q)\) must be listed in type \(\omega\)
\cite[\S2, p.~16]{Keisler1970}.
\end{hbexample}

\begin{hbwarning}{def:p02-w02c-p02-003-kq-axiom-zero}{The first-order part of the axiom system}
Keisler's Axiom~0 is ``all the axiom schemes for \(L\)'' with a reference
to Church's textbook \cite[\S1, p.~6]{Keisler1970}; that list was not
compared here.  The proof in the annex uses Axiom~0 only through
(0a)--(0g): soundness needs that each scheme holds in every weak model, and
the Henkin stage needs the deduction theorem, the derived rules and the
lemma on new constants of \cref{thm:p02-w02c-p02-003-kq-proof-lemmas}.  Any
list of first-order schemes for which these lemmas are proved yields the
theorems with the same proof.
\end{hbwarning}

\begin{proofboundary}{def:p02-w02c-p02-003-beyond-boundary}{Claims of Barwise's \S5 whose proofs are not written in the handbook}
\begin{enumerate}
\item \emph{The least logic containing weak second-order logic in which
      implicit definitions become explicit} (p.~44).  The chapter says only
      that the search for such a logic leads to the admissible fragments
      \(L_A\) of \(L_{\omega_1\omega}\) of Chapter~A.7, and gives no precise
      statement.  The precise result in Barwise's paper is his
      Theorem~4.1: for the least logic \(\Delta(\mathcal L^{\ast})\) that is
      as strong as \(\mathcal L^{\ast}\) and has the Souslin--Kleene
      property (every class that is \(\Delta^1_1\) in the logic is
      definable in it), \(\Delta(L^{w})\), \(\Delta(L(Q_0))\) and
      \(L_{\mathrm{HYP}}\) are equivalent, where \(L^{w}\) is weak
      second-order logic, \(\mathrm{HYP}\) is the least admissible set with
      \(\omega\) as an element, and \(L_{\mathrm{HYP}}\) is its admissible
      fragment restricted to formulas with finitely many symbols
      \cite[Part~II, \S4, 4.1, p.~260]{Barwise1974Axioms}.  None of the
      following is written in the handbook: the construction of
      \(\Delta(\mathcal L^{\ast})\) as a logic in the sense of
      \cref{def:p02-w02c-p02-003-b74-logic}, whose details the paper refers
      to Makowsky, Shelah and Stavi; the interpolation property of
      \(L_{\mathrm{HYP}}\), which the paper takes from Barwise's earlier work;
      the three facts about admissible sets with urelements that the proof
      quotes from Barwise's book then in preparation
      \cite[Part~II, \S4, pp.~260--262]{Barwise1974Axioms}; and the passage
      from the Souslin--Kleene property, which concerns
      \(\Delta^1_1\) classes, to the chapter's wording in terms of implicit
      and explicit definability of relations.  What the handbook supplies
      toward the claim is the failure of Beth's theorem for weak
      second-order logic (\cref{thm:p02-w02c-p02-003-wso-beth-failure}),
      Beth's theorem for \(L_{\omega_1\omega}\)
      (\cref{thm:p02-w02c-p02-003-beth}), the translation of weak
      second-order logic into \(L_{\omega_1\omega}\)
      (\cref{thm:p02-w02c-p02-003-wso-translation}) and the admissibility of
      \(\mathrm{HF}(\mathfrak M)\) (\cref{thm:p02-w02c-p02-003-hf}).
\end{enumerate}
\end{proofboundary}

\subsection{Barwise's claims about particular theories}
\label{sec:p02-w02c-p02-003-theories}

\begin{hbdefinition}{def:p02-w02c-p02-003-decidable}{Complete and decidable theories}
A set \(T\) of sentences of \(L\) is \emph{complete} if for every sentence
\(\sigma\) of \(L\), \(T\models\sigma\) or \(T\models\neg\sigma\).  The
theory axiomatized by \(T\) is \emph{decidable} if there is an effective
procedure that tells of an arbitrary sentence \(\sigma\) of \(L\) whether
\(T\models\sigma\); a structure \(\mathfrak A\) is decidable if its theory
\(\{\sigma:\mathfrak A\models\sigma\}\) is
\cite[ch.~A.1, \S2, p.~16]{Barwise1977FirstOrder}.  Here ``effective
procedure'' is used in Barwise's informal sense.  For \(p=0\) or a prime,
\(\mathrm{ACF}_p\) is the set of sentences of \(\{+,\cdot,0,1\}\)
consisting of the field axioms of \cref{def:p02-w02-p02-005-bw-rings}, for
each \(n\geq1\) the sentence
\(\forall y_0\cdots\forall y_{n-1}\,\exists x\,
[x^n+y_{n-1}x^{n-1}+\cdots+y_0=0]\), and \(1+\cdots+1=0\) with \(p\)
summands if \(p\) is prime, or \(\neg(1+\cdots+1=0)\) with \(k\) summands for
every \(k\geq1\) if \(p=0\).
\end{hbdefinition}

\begin{hblemma}{thm:p02-w02c-p02-003-los-vaught}{proved=annex:p02-w02c-p02-003-los-vaught}{Categoricity test for completeness}
Let \(T\) be a set of first-order sentences of \(L\) that has no finite
model, and is \(\kappa\)-categorical (any two models of
cardinality \(\kappa\) are isomorphic) for some cardinal
\(\kappa\geq\max(\abs{L},\aleph_0)\).  Then \(T\) is complete.  The axiom
of choice is assumed.
\end{hblemma}

Barwise refers the results about particular theories to the chapters on
recursion theory, decidable theories and incompleteness of the same
handbook (p.~16).  Rabin's chapter on decidable theories proves the
categoricity test for countable theories
\cite[\S2, 2.1, Theorem~4, p.~608]{Rabin1977Decidable}
and derives the completeness of each \(\mathrm{ACF}_p\) from Steinitz's
structure theorem for algebraically closed fields, which it uses without
proof \cite[\S2, 2.2, p.~609]{Rabin1977Decidable}.  The
algebra needed for that structure theorem, and the fundamental theorem of
algebra, are proved next from the field axioms, the axiom of choice (in
the form of Zorn's lemma) and, for the fundamental theorem, the
completeness of the real numbers.  For fields \(K\subseteq E\), an element
\(a\in E\) is \emph{algebraic over \(K\)} if \(q(a)=0\) for some nonzero
\(q\in K[x]\); \(K[a]=\{q(a):q\in K[x]\}\).  A subset \(X\) of a field
\(E\) is \emph{algebraically independent} over a subfield \(P\) if
\(q(x_1,\ldots,x_k)\neq0\) for all distinct \(x_1,\ldots,x_k\in X\) and
every nonzero polynomial \(q\in P[t_1,\ldots,t_k]\); a \emph{transcendence
basis} of \(E\) over \(P\) is a maximal such set.  \(P(X)\) is the least
subfield of \(E\) containing \(P\cup X\).  The \emph{prime field} of
\(E\) is its least subfield.  A field \(E\) is \emph{algebraically
closed} if every polynomial \(x^n+y_{n-1}x^{n-1}+\cdots+y_0\) with
\(n\geq1\) and coefficients in \(E\) has a root in \(E\), that is, if
\(E\) satisfies the axioms of \(\mathrm{ACF}_p\) other than the
characteristic axioms.

\begin{hblemma}{thm:p02-w02c-p02-003-simple-extension}{proved=annex:p02-w02c-p02-003-simple-extension}{Simple algebraic extensions}
Let \(K\subseteq E\) be fields and let \(a\in E\) be algebraic over \(K\).
\begin{enumerate}
\item There is a monic irreducible \(m\in K[x]\) with \(m(a)=0\) that
      divides every \(q\in K[x]\) with \(q(a)=0\).
\item \(K[a]\) is a subfield of \(E\).
\item If \(G\) is an algebraically closed field and \(\sigma\colon K\to G\)
      is an embedding of fields, then \(\sigma\) extends to an embedding
      \(K[a]\to G\).
\item If \(E\) is algebraically closed, every \(b\) algebraic over \(E\) in
      a field extending \(E\) lies in \(E\).
\end{enumerate}
\end{hblemma}

\begin{hbtheorem}{thm:p02-w02c-p02-003-steinitz}{proved=annex:p02-w02c-p02-003-steinitz}{Steinitz's classification in uncountable cardinalities}
Assume the axiom of choice.  Let \(E\) and \(G\) be algebraically closed
fields of the same characteristic \(p\) (\(0\) or a prime).
\begin{enumerate}
\item \(E\) has a transcendence basis \(X\) over its prime field \(P\), and
      every element of \(E\) is algebraic over \(P(X)\).
\item If \(E\) is uncountable, every such \(X\) has cardinality
      \(\abs{E}\).
\item If \(E\) and \(G\) have transcendence bases of the same cardinality
      over their prime fields, then \(E\cong G\).  In particular, if
      \(\abs{E}=\abs{G}\) is uncountable, then \(E\cong G\): for every
      uncountable cardinal \(\kappa\), \(\mathrm{ACF}_p\) is
      \(\kappa\)-categorical.
\end{enumerate}
\end{hbtheorem}

\begin{hbtheorem}{thm:p02-w02c-p02-003-fta}{proved=annex:p02-w02c-p02-003-fta}{Fundamental theorem of algebra}
The field \(\mathbb C\) of complex numbers is algebraically closed.  Hence
\(\mathbb C\) is a model of \(\mathrm{ACF}_0\); in particular an
algebraically closed field of characteristic \(0\) exists.
\end{hbtheorem}

\begin{hbtheorem}{thm:p02-w02c-p02-003-acf}{proved=annex:p02-w02c-p02-003-acf}{Algebraically closed fields}
Assume the axiom of choice.
\begin{enumerate}
\item For \(p=0\) or a prime, \(\mathrm{ACF}_p\) is complete: any two
      algebraically closed fields of the same characteristic satisfy the
      same first-order sentences of \(\{+,\cdot,0,1\}\).
\item There is an effective procedure that decides, for a sentence
      \(\sigma\) of \(\{+,\cdot,0,1\}\), whether
      \(\mathrm{ACF}_0\models\sigma\).
\item For the field \(\mathbb C\) of complex numbers,
      \(\mathbb C\models\sigma\) iff \(\mathrm{ACF}_0\models\sigma\), and
      the procedure of part~2 decides the theory of \(\mathbb C\).
\end{enumerate}
Part~1 rests on the uncountable categoricity of \(\mathrm{ACF}_p\),
\cref{thm:p02-w02c-p02-003-steinitz} (stated with a cited proof as
\cref{thm:p03-acfp-uncountable-categorical}); parts~2 and~3 rest on the
fundamental theorem of algebra, \cref{thm:p02-w02c-p02-003-fta}.
\end{hbtheorem}

\begin{hbdefinition}{def:p02-w02c-p02-003-ring-coding}{Translation into the ring language and its coding}
Formulas of \(\{+,\cdot,0,1\}\) are built with the variables \(x_i\) and the
connectives and quantifiers of \(\Lang_A\)
(\cref{sec:arithmetic-coding-04}).  For such a formula \(\varphi\), let
\(\varphi^{\circ}\) be the \(\Lang_A\)-formula obtained by writing
\(S(0)\) for \(1\) and \(\times\) for \(\cdot\), and put
\(\ulcorner\varphi\urcorner'=\ulcorner\varphi^{\circ}\urcorner\), the
G\"odel number of \cref{sec:arithmetic-coding-04}.  For each \(i\), let
\(D(x_i)\) be the formula
\[
 \exists x_{i+1}\exists x_{i+2}\exists x_{i+3}\exists x_{i+4}\,
 \bigl(x_i=((x_{i+1}\cdot x_{i+1}+x_{i+2}\cdot x_{i+2})
 +x_{i+3}\cdot x_{i+3})+x_{i+4}\cdot x_{i+4}\bigr).
\]
For an \(\Lang_A\)-term \(t\), \(t^{\mathbb Z}\) is the term of
\(\{+,\cdot,0,1\}\) with \(x_i^{\mathbb Z}=x_i\), \(0^{\mathbb Z}=0\),
\((S\,t)^{\mathbb Z}=t^{\mathbb Z}+1\),
\((t+t')^{\mathbb Z}=t^{\mathbb Z}+t'^{\mathbb Z}\) and
\((t\times t')^{\mathbb Z}=t^{\mathbb Z}\cdot t'^{\mathbb Z}\).  For an
\(\Lang_A\)-formula \(\chi\), \(\chi^{\mathbb Z}\) applies this to the
terms of atomic formulas, commutes with \(\bot\) and the connectives, and
sends \(\forall x_i\,\chi\) to \(\forall x_i\,(D(x_i)\to\chi^{\mathbb
Z})\) and \(\exists x_i\,\chi\) to \(\exists x_i\,(D(x_i)\land\chi^{\mathbb
Z})\).
\end{hbdefinition}

The only free variable of \(D(x_i)\) is \(x_i\), so the choice of the
bound variables \(x_{i+1},\ldots,x_{i+4}\) cannot capture a variable of
\(\chi\).  The map \(\varphi\mapsto\varphi^{\circ}\) is injective, since
\(\varphi\) is recovered from \(\varphi^{\circ}\) by writing \(1\) for
\(S(0)\) and \(\cdot\) for \(\times\); so
\(\ulcorner\cdot\urcorner'\) is a coding of the formulas of
\(\{+,\cdot,0,1\}\).

\begin{hbtheorem}{thm:p02-w02c-p02-003-lagrange}{proved=annex:p02-w02c-p02-003-lagrange}{Lagrange's four-square theorem}
Every natural number is a sum of four squares of integers.  Hence the
formula \(D(x)\) of \cref{def:p02-w02c-p02-003-ring-coding} defines the
set of non-negative integers in \(\mathbb Z\).
\end{hbtheorem}

\begin{hblemma}{thm:p02-w02c-p02-003-ring-coding}{proved=annex:p02-w02c-p02-003-ring-coding}{The translation is arithmetically definable}
There is a primitive recursive function \(\mathrm{tr}\) with
\(\mathrm{tr}(\ulcorner\chi\urcorner)=\ulcorner\chi^{\mathbb
Z}\urcorner'\) for every \(\Lang_A\)-formula \(\chi\).  The relation
\(\{(\ulcorner\sigma\urcorner,\ulcorner\sigma^{\mathbb
Z}\urcorner'):\sigma\text{ a first-order }\Lang_A\text{-sentence}\}\) is
defined in \(\mathcal N\) by a first-order formula of \(\Lang_A\).
\end{hblemma}

\begin{hbtheorem}{thm:p02-w02c-p02-003-integers}{proved=annex:p02-w02c-p02-003-integers}{The ring of integers is not decidable (Barwise, p.~16)}
Let \(\mathbb Z=(\mathbb Z,+,\cdot,0,1)\), let \(\mathcal N\) and
\(\mathrm{TA}\) be as in \cref{def:p02-w02-p02-040-so-coding}.  For a
first-order sentence \(\sigma\) of \(\Lang_A=\{0,S,+,\times\}\), let
\(\sigma^{\mathbb Z}\) be the sentence of \(\{+,\cdot,0,1\}\) of
\cref{def:p02-w02c-p02-003-ring-coding}, obtained by writing \(t+1\) for
\(S\,t\) and \(\cdot\) for \(\times\) and relativizing every quantifier to
\(D\).
\begin{enumerate}
\item \(\mathcal N\models\sigma\) iff
      \(\mathbb Z\models\sigma^{\mathbb Z}\) for every \(\sigma\).
\item Let \(\ulcorner\cdot\urcorner'\) be the coding of
      \cref{def:p02-w02c-p02-003-ring-coding}, or any coding of the
      sentences of \(\{+,\cdot,0,1\}\) by natural numbers for which the
      relation
      \(\{(\ulcorner\sigma\urcorner,\ulcorner\sigma^{\mathbb Z}\urcorner'):
      \sigma\text{ a first-order }\Lang_A\text{-sentence}\}\) is definable
      in \(\mathcal N\) by a first-order formula (the first is one by
      \cref{thm:p02-w02c-p02-003-ring-coding}).  Then the set
      \(W=\{\ulcorner\varphi\urcorner':\mathbb Z\models\varphi\}\) is not
      definable in \(\mathcal N\) by any first-order formula; in particular
      it is not computably enumerable, and so no effective procedure
      decides the theory of \(\mathbb Z\).
\item For such a coding, no set \(T\) of sentences true in
      \(\mathbb Z\) whose set of consequences
      \(\{\ulcorner\varphi\urcorner':T\models\varphi\}\) is computably
      enumerable is complete.
\end{enumerate}
\end{hbtheorem}

Part~3 shows that some true sentence is missing from the consequences of
\(T\) without naming one.  Barwise adds that G\"odel's second
incompleteness theorem shows how to find a sentence true in \(\mathbb Z\)
that does not follow from \(T\)
\cite[ch.~A.1, \S2, p.~16]{Barwise1977FirstOrder}.  The rest of this
subsection names that sentence: the translation of the consistency sentence
of an arithmetic theory \(T^{\sharp}\) read off from \(T\), formed with a
provability formula selected so that the derivability conditions can be
proved for it.  The construction uses \(\mathsf Q\) and \(\PA\)
(\cref{sec:arithmetic-coding-01}); the fixed G\"odel numbering, the
twenty-rule derivation records and the certified axiom leaves of
\cref{sec:arithmetic-coding-04}; the literal presentation \(\PA^{\rm d}\)
of \(\PA\) with its recognizer \(\mathsf{Ax}^{\rm d}_{\PA}\)
(\cref{def:p05-direct-pa-axioms}); the selected formula
\(\mathsf{Prv}^{19}_{\PA}\), its characteristic description
\(\chi_{19}\) and its axiom description \(\mathsf{Ax}_{19}\)
(\cref{def:p05-selected-direct-pa-graph,def:p05-checker-axiom-description});
the graph formulas \(\Gamma_d\) of \cref{def:p05-pa-pr-graph}; and the
derivability conditions D1--D3 of \cref{sec:second-incompleteness-01}.
\(\mathcal N\) is the standard structure \((\N,0,S,+,\times)\) of
\cref{def:language-arithmetic}, as in
\cref{thm:p02-w02c-p02-003-integers}.

Barwise speaks of a list of true axioms given by an effective procedure.
Read through the Church--Turing thesis (\cref{def:p04-church-turing}), this
is the hypothesis that the set of codes of the list is computably
enumerable.  The results below assume exactly that hypothesis; the thesis is
used in none of their proofs.

\begin{hbdefinition}{def:p02-w02c-p02-003-sharp-enumeration}{The theory read off from a computably enumerable theory of \(\mathbb Z\)}
Let \(T\) be a set of sentences of \(\{+,\cdot,0,1\}\) for which the set
\(C_T=\{\ulcorner\tau\urcorner':\tau\in T\}\) is computably enumerable
(\cref{sec:turing-machines-03}), with \(\ulcorner\cdot\urcorner'\) and
\(\tau^{\circ}\) as in \cref{def:p02-w02c-p02-003-ring-coding}.  Fix an
index \(e_T\) whose partial computable function has domain \(C_T\)
(\cref{thm:enumeration}), and let \(V(e,a,s)\) be the primitive recursive
truth value of the bounded formula
\(\mathsf{Comp}(\bar e,\bar a,\bar s)\) fixed in
\cref{sec:arithmetic-coding-04}.  Put
\[
 T^{+}=T\cup\{\alpha^{\mathbb Z}:\alpha\in\PA^{\rm d}\},\qquad
 T^{\sharp}=\{\sigma:\sigma\text{ an }\Lang_A\text{-sentence and }
 T^{+}\models\sigma^{\mathbb Z}\},
\]
and write \((T^{+})^{\circ}=\{\chi^{\circ}:\chi\in T^{+}\}\).  With
\(\mathrm{tr}\) the primitive recursive function of
\cref{thm:p02-w02c-p02-003-ring-coding}, define the leaf test
\[
 \mathsf{Lf}(a,s)=\mathsf{sg}\bigl(V(e_T,a,s)
       +E(a,\mathrm{tr}(s))\cdot\mathsf{Ax}^{\rm d}_{\PA}(s)\bigr),
\]
where \(\mathsf{Ax}^{\rm d}_{\PA}\) also names its characteristic
function, \(E\) is the equality test and \(\mathsf{sg}\) the sign
(\cref{def:p05-checker-description-language}).  Let
\(\mathsf{Pf}^{+}(p,y)\) hold iff \(p=N_{16}(N_{15}(R),W)\), the row list
\(R\) passes all the derivation checks of \cref{sec:arithmetic-coding-04},
its root formula is a sentence with code \(y\), and, for the root open list
\([j_0,\ldots,j_{u-1}]\), \(W=[s_0,\ldots,s_{u-1}]\) (with \(W=0\) when
\(u=0\)) and every leaf formula code \(a_{j_v}\) is a sentence code with
\(\mathsf{Lf}(a_{j_v},s_v)=1\).  This is the certified relation
\(\mathsf{Proof}_T\) of \cref{sec:arithmetic-coding-04} with the leaf test
\(V(e_T,a_{j_v},s_v)=1\) replaced by \(\mathsf{Lf}(a_{j_v},s_v)=1\).
Finally let
\[
 \mathsf{En}^{\sharp}(p,x)\iff x\text{ is the code of an }
 \Lang_A\text{-sentence and }\mathsf{Pf}^{+}(p,\mathrm{tr}(x)).
\]
\end{hbdefinition}

A certificate \(s_v\) is thus either an accepting computation showing that
the leaf belongs to \(T^{\circ}\), or the code of an axiom
\(\alpha\in\PA^{\rm d}\) whose translation is the leaf.  The relation
\(\mathsf{En}^{\sharp}\) is the explicit enumeration of \(T^{\sharp}\): a
number \(p\) with \(\mathsf{En}^{\sharp}(p,\ulcorner\sigma\urcorner)\) is
a certified derivation of \((\sigma^{\mathbb Z})^{\circ}\).

\begin{hblemma}{thm:p02-w02c-p02-003-ring-transfer}{proved=annex:p02-w02c-p02-003-ring-transfer}{Consequence in the ring language and in \(\Lang_A\)}
For every set \(\Sigma\) of sentences of \(\{+,\cdot,0,1\}\) and every
sentence \(\varphi\) of that language,
\(\Sigma\models\varphi\) iff
\(\{\psi^{\circ}:\psi\in\Sigma\}\models\varphi^{\circ}\).
\end{hblemma}

\begin{hblemma}{thm:p02-w02c-p02-003-sharp-enumeration}{proved=annex:p02-w02c-p02-003-sharp-enumeration}{An explicit enumeration of \(T^{\sharp}\)}
Let \(T\) be as in \cref{def:p02-w02c-p02-003-sharp-enumeration}.
\begin{enumerate}
\item \(\mathsf{Pf}^{+}\) and \(\mathsf{En}^{\sharp}\) are primitive
      recursive.
\item For every \(\Lang_A\)-sentence \(\psi\),
      \(\exists p\,\mathsf{Pf}^{+}(p,\ulcorner\psi\urcorner)\) iff
      \((T^{+})^{\circ}\derives\psi\).
\item The set of codes of members of \(T^{\sharp}\) is
      \(\{x:\exists p\,\mathsf{En}^{\sharp}(p,x)\}\).  It is the domain
      of the partial recursive function
      \(x\mapsto\mu p\,[Z(\chi_{\mathsf{En}^{\sharp}}(p,x))=0]\), where
      \(\chi_{\mathsf{En}^{\sharp}}\) is the characteristic function of
      \(\mathsf{En}^{\sharp}\); hence \(T^{\sharp}\) is computably
      enumerable.
\end{enumerate}
\end{hblemma}

Part~3 passes from the partial recursive function to a Turing machine by
\cref{thm:p04-tm-partial-recursive}, which the handbook cites
(\cref{def:p02-w02c-p02-003-theories-boundary}, item~3).  The axiom set
\(T^{\sharp}\) is only computably enumerable, while the factor
\(\mathsf{Ax}_{19}\) of the selected \(\PA\) formula is a primitive
recursive test on single leaf formulas.  Craig's device of padding each
member of \(T^{\sharp}\) by its own enumeration witness gives a primitive
recursive axiom set with the same consequences, and the selected formula
for \(T^{\sharp}\) adds a test for it to \(\mathsf{Ax}_{19}\).

\begin{hbdefinition}{def:p02-w02c-p02-003-sharp-formula}{A selected provability formula for \(T^{\sharp}\)}
Let \(T\) be as in \cref{def:p02-w02c-p02-003-sharp-enumeration}.  Put
\[
 \mathrm{pd}(0,c)=c,\qquad
 \mathrm{pd}(k+1,c)=N_8(N_5(N_1(),N_1()),\mathrm{pd}(k,c)),
\]
so that
\(\mathrm{pd}(k,\ulcorner\sigma\urcorner)\) is the code of
\(0=0\land(0=0\land\cdots\land(0=0\land\sigma)\cdots)\) with \(k\)
conjuncts \(0=0\).  Let \(\mathcal A^{\sharp}\) be the set of sentences
\(\theta\) with \(\ulcorner\theta\urcorner=\mathrm{pd}(k,x)\) for some
\(k,x\) such that \(\mathsf{En}^{\sharp}(k,x)\), and put
\[
 \mathrm{cr}(a)=\mathsf{sg}\Bigl(\sum_{k\le a}\sum_{x\le a}
     E(a,\mathrm{pd}(k,x))\,\chi_{\mathsf{En}^{\sharp}}(k,x)\Bigr).
\]
Fix a primitive recursive description \(\mathsf{Cr}\) of \(\mathrm{cr}\)
in the sense of \cref{def:p05-pa-pr-graph}, and put
\[
 \mathsf{Ax}_{\sharp}(a)=\mathsf{sg}\bigl(A(\mathsf{Ax}_{19}(a),
 \mathsf{Cr}(a))\bigr),
\]
with the library descriptions \(\mathsf{sg},A\) of
\cref{def:p05-checker-description-language}.  Let \(\chi_{\sharp}\) be the
description \(\chi_{19}\) of \cref{def:p05-selected-direct-pa-graph} with
each factor \(\mathsf{Ax}_{19}(f_R(e(\mathsf{open}(R),v)))\) of its final
product replaced by \(\mathsf{Ax}_{\sharp}(f_R(e(\mathsf{open}(R),v)))\)
and every other part unchanged, and let \(d_{\sharp}\) be its expansion.
Select
\[
 \mathsf{Prf}^{\sharp}(p,x)\defeq\Gamma_{d_{\sharp}}(p,x,\bar1),\qquad
 \mathsf{Prv}^{\sharp}(x)\defeq\exists p\,\mathsf{Prf}^{\sharp}(p,x),
\]
\[
 \mathsf{Con}^{\sharp}\defeq\neg\mathsf{Prv}^{\sharp}
       (\overline{\ulcorner\bottom_A\urcorner}),
\]
with the variable conventions of \(\mathsf{Prv}^{19}_{\PA}\).  Let
\(\mathsf{mp}_{\sharp}\) and its expanded description
\(d^{\sharp}_{\rm mp}\) be \(\mathsf{mp}_{19}\) and \(d_{\rm mp}\) of
\cref{def:p05-direct-mp-description} with \(\chi_{\sharp}\) in place of
\(\chi_{19}\) in the success guard.
\end{hbdefinition}

The formula \(\mathsf{Prv}^{\sharp}\) depends on \(T\) and on the choices of
\(e_T\) and \(\mathsf{Cr}\); the results below hold for every choice.  With
\(\mathsf{Prov}_{T^{\sharp}}\) taken to be \(\mathsf{Prv}^{\sharp}\), as
\cref{sec:second-incompleteness-01} permits for a specified formula,
\(\mathsf{Con}^{\sharp}\) is the sentence \(\mathsf{Con}(T^{\sharp})\) of
\cref{def:formal-consistency-sentence}.

\begin{hbtheorem}{thm:p02-w02c-p02-003-sharp-provability}{proved=annex:p02-w02c-p02-003-sharp-provability}{Proof existence and the derivability conditions for \(\mathsf{Prv}^{\sharp}\)}
Let \(T\) be as in \cref{def:p02-w02c-p02-003-sharp-enumeration}.
\begin{enumerate}
\item \(\PA^{\rm d}\subseteq T^{\sharp}\), and every theorem of \(\PA\)
      is a theorem of \(T^{\sharp}\).  The formulas \(\mathsf{Prf}^{\sharp}\)
      and \(\mathsf{Prv}^{\sharp}\) are \(\Sigma^0_1\), and for every
      \(\Lang_A\)-sentence \(\sigma\),
      \[
       T^{\sharp}\derives\sigma
       \quad\Longleftrightarrow\quad
       \PA^{\rm d}\cup\mathcal A^{\sharp}\derives\sigma
       \quad\Longleftrightarrow\quad
       \mathcal N\models\mathsf{Prv}^{\sharp}
                        (\overline{\ulcorner\sigma\urcorner}).
      \]
\item (D1) If \(T^{\sharp}\derives\sigma\), then
      \(\PA\derives\mathsf{Prv}^{\sharp}(\overline{\ulcorner\sigma\urcorner})\).
\item In \(\PA\), with universal closure in \(p,q,u,v\),
      \[
       \mathsf{Prf}^{\sharp}(p,N_{10}(u,v))\land\mathsf{Prf}^{\sharp}(q,u)
       \longrightarrow\exists r\,\bigl(\Gamma_{d^{\sharp}_{\rm mp}}(p,q,r)
       \land\mathsf{Prf}^{\sharp}(r,v)\bigr),
      \]
      the constructor being eliminated through its specified graph.
      Hence (D2) for all sentences \(\varphi,\psi\),
      \(\PA\derives\mathsf{Prv}^{\sharp}
      (\overline{\ulcorner\varphi\to\psi\urcorner})\to
      (\mathsf{Prv}^{\sharp}(\overline{\ulcorner\varphi\urcorner})\to
      \mathsf{Prv}^{\sharp}(\overline{\ulcorner\psi\urcorner}))\).
\item \(\PA\derives\forall p\,\forall x\,
      (\mathsf{Prf}^{19}_{\PA}(p,x)\to\mathsf{Prf}^{\sharp}(p,x))\).
\item (D3) For every sentence \(\varphi\),
      \(\PA\derives\mathsf{Prv}^{\sharp}(\overline{\ulcorner\varphi\urcorner})
      \to\mathsf{Prv}^{\sharp}(\overline{\ulcorner\mathsf{Prv}^{\sharp}
      (\overline{\ulcorner\varphi\urcorner})\urcorner})\).
\end{enumerate}
By part~1, the conclusions of parts~2, 3 and~5 hold with \(T^{\sharp}\) in
place of \(\PA\); so \(\mathsf{Prv}^{\sharp}\) satisfies D1--D3 for the
theory \(T^{\sharp}\).
\end{hbtheorem}

The internal arguments for parts~3--5 are those of
\cref{thm:p05-direct-pa-internal-mp,thm:p05-direct-pa-d3} for
\(\mathsf{Prv}^{19}_{\PA}\): the joining argument uses the axiom
description only through the uniqueness of its graph outputs, and D3 for
\(\mathsf{Prv}^{\sharp}\) is obtained from the \(\Sigma^0_1\) introspection
of \(\PA\) followed by part~4.

\begin{hbtheorem}{thm:p02-w02c-p02-003-integers-second}{proved=annex:p02-w02c-p02-003-integers-second}{A sentence true in \(\mathbb Z\) that does not follow from a given effective true theory (Barwise, p.~16)}
Let \(\sigma\mapsto\sigma^{\mathbb Z}\) be the translation of
\cref{thm:p02-w02c-p02-003-integers}.  Let \(T\) be a set of sentences of
\(\{+,\cdot,0,1\}\), each true in \(\mathbb Z\), whose set of codes is
computably enumerable, and let \(T^{+}\), \(T^{\sharp}\) and
\(\mathsf{Con}^{\sharp}\) be as in
\cref{def:p02-w02c-p02-003-sharp-enumeration,def:p02-w02c-p02-003-sharp-formula}.
\begin{enumerate}
\item \(T^{\sharp}\supseteq\PA^{\rm d}\); in particular \(T^{\sharp}\)
      contains the seven axioms of \(\mathsf Q\) in the literal form of
      \cref{sec:p05-standard-pa-checker}.  Every member of \(T^{\sharp}\) is
      true in \(\mathcal N\), \(T^{\sharp}\) is consistent, and \(T^{\sharp}\)
      is computably enumerable.
\item The sentence \((\mathsf{Con}^{\sharp})^{\mathbb Z}\) of
      \(\{+,\cdot,0,1\}\) is true in \(\mathbb Z\), and
      \(T^{+}\not\models(\mathsf{Con}^{\sharp})^{\mathbb Z}\); in
      particular \(T\not\models(\mathsf{Con}^{\sharp})^{\mathbb Z}\).
\end{enumerate}
\end{hbtheorem}

The sentence \((\mathsf{Con}^{\sharp})^{\mathbb Z}\) is built from \(T\) in
a fixed way: an index \(e_T\) for \(C_T\) gives the relations
\(\mathsf{En}^{\sharp}\) and \(\mathrm{cr}\), a description of
\(\mathrm{cr}\) gives \(d_{\sharp}\), hence \(\mathsf{Prv}^{\sharp}\),
\(\mathsf{Con}^{\sharp}\) and its translation.  This is the sense in which
the second incompleteness theorem shows how to find a true sentence that
\(T\) leaves unproved.

\emph{Worked example.}  Let \(T=\varnothing\).  Then \(C_T=\varnothing\),
\(T^{+}=\{\alpha^{\mathbb Z}:\alpha\in\PA^{\rm d}\}\), and
\(T^{\sharp}\) is the set of \(\Lang_A\)-sentences \(\sigma\) whose
translations follow from the translated axioms of \(\PA\).  For \(e_T\) any
index of the everywhere undefined function, \(V(e_T,a,s)=0\) for all
\(a,s\), so a certificate in \(\mathsf{Pf}^{+}\) is always the code of an
axiom of \(\PA^{\rm d}\).  The theorem yields a sentence of
\(\{+,\cdot,0,1\}\), true in \(\mathbb Z\), that does not follow from the
translated axioms of \(\PA\).  Its negation does not follow from them
either, since \(\mathbb Z\) is a model of \(T^{+}\) in which the sentence
is true.  So, although \(\PA\) is interpreted in \(\mathbb Z\) through the
four-square formula \(D\), the translated axioms leave this translated
consistency statement undecided.

The selection of the provability formula cannot be dropped.  The following
lemma shows that for an arbitrary \(\Sigma^0_1\) formula numeralwise
representing the certified proof relation, as allowed in
\cref{def:fixed-proof-predicate}, the conclusion of the second
incompleteness theorem can fail; so D2 or D3 can fail for such a formula,
and no proof of them for the formula of \cref{def:fixed-proof-predicate}
alone is possible.

\begin{hblemma}{thm:p02-w02c-p02-003-arbitrary-representative}{proved=annex:p02-w02c-p02-003-arbitrary-representative}{A numeralwise representative with a provable consistency sentence}
Let \(S\supseteq\mathsf Q\) be a consistent computably enumerable theory,
let \(\mathsf{Proof}_S\) be its certified proof relation
(\cref{sec:arithmetic-coding-04}), and let \(\rho(p,x)\) be a
\(\Sigma^0_1\) formula that numeralwise represents \(\mathsf{Proof}_S\) in
\(\mathsf Q\).  Put \(c=\ulcorner\bottom_A\urcorner\) and
\(\rho'(p,x)\defeq\rho(p,x)\land\neg\,x=\bar c\).  Then \(\rho'\) is a
\(\Sigma^0_1\) formula that numeralwise represents \(\mathsf{Proof}_S\) in
\(\mathsf Q\); \(S\derives\sigma\) iff
\(\mathcal N\models\exists p\,\rho'(p,\overline{\ulcorner\sigma\urcorner})\)
for every sentence \(\sigma\); \(\exists p\,\rho'\) satisfies D1 for \(S\);
and \(S\derives\neg\exists p\,\rho'(p,\bar c)\).  Consequently
\(\exists p\,\rho'\) does not satisfy both D2 and D3 for \(S\).
\end{hblemma}

\begin{proofboundary}{def:p02-w02c-p02-003-theories-boundary}{Claims of Barwise's p.~16 whose proofs are not written in the handbook}
\begin{enumerate}
\item \emph{The theory of abelian groups is decidable} (Szmielew's
      theorem).  Missing: an effective reduction of every sentence of
      \(\{+,0\}\) to a Boolean combination of invariant sentences about
      the numbers of elements of order \(p\), divisibility and torsion, and
      a proof that the reduction preserves truth in every abelian group.
      Rabin's chapter on decidable theories states the theorem without
      proof \cite[p.~611, Theorem~7]{Rabin1977Decidable}.
\item \emph{The theory of groups is undecidable}.  Missing: an
      interpretation, in a model class of a finitely axiomatized group
      theory, of an essentially undecidable theory, and the transfer of
      undecidability along it (the transfer itself is the cited
      \cref{thm:p05-interpretation-transfer}).
\item In \cref{thm:p02-w02c-p02-003-sharp-enumeration}, part~3: the
      passage from the explicit partial recursive function whose domain is
      the set of codes of \(T^{\sharp}\) to a Turing machine with that
      domain.  It uses the cited \cref{thm:p04-tm-partial-recursive}
      (partial recursive functions are Turing computable); computable
      enumerability of \(T^{\sharp}\) in the machine sense of
      \cref{sec:turing-machines-03}, and through it the hypothesis of
      \cref{thm:godel-second} used in
      \cref{thm:p02-w02c-p02-003-integers-second}, rests on that cited
      theorem.  The enumeration itself, the derivability conditions for
      \(\mathsf{Prv}^{\sharp}\) and the rest of the argument are proved.
\end{enumerate}
Barwise refers these results to Chapters~C.1, C.3 and D.1 of the same
handbook \cite[ch.~A.1, \S2, p.~16]{Barwise1977FirstOrder}.
\end{proofboundary}

\emph{Relationships.}  \Cref{thm:p02-w02c-p02-003-general-lindstrom}
reduces to \cref{thm:p02-w02b-p02-001-barwise-lindstrom} and through it
to \cref{thm:p03-lindstrom}.  \Cref{thm:p02-w02c-p02-003-b74-lindstrom}
does not reduce to either: it adapts the coding of the proof of
\cref{thm:p03-lindstrom} to function symbols and to Barwise's translation
axiom, and adds the indexed-union argument for the upward
L\"owenheim--Skolem hypothesis.  \Cref{thm:p02-w02c-p02-003-or-not} rests on
\cref{thm:p02-w02-p02-008-neg-and}, and
\cref{thm:p02-w02c-p02-003-small-model} on \cref{thm:fol-completeness} and
\cref{thm:downward-ls}.
\Cref{thm:p02-w02c-p02-003-infinitary-rule} is a consequence of the
failure of compactness in part~3 of
\cref{thm:p02-w02-p02-005-bw-extended-ls}.
\Cref{thm:p02-w02c-p02-003-infinitary-completeness} is its converse for the
calculus \(\mathbf H_{\infty}\); \(\mathbf H_{\infty}\) carries Barwise's H
(\cref{def:p02-w02-p02-005-bw-system-h}) to sentences of
\(L_{\omega_1\omega}\) with parameters, and its completeness proof replaces
Barwise's reduction to propositional compactness, which is not available
for \(L_{\omega_1\omega}\), by
\cref{thm:p02-w02c-p02-003-term-model-existence}, the extension of
\cref{thm:p02-w02c-p02-003-model-existence} to function symbols.  It is a
different calculus from the one of the cited
\cref{thm:p02-karp-completeness}.
\Cref{thm:p02-w02c-p02-003-wso-translation} complements part~2 of that
theorem, which translates \(L(Q_0)\) into \(L_{\omega_1\omega}\).
\Cref{thm:p02-w02c-p02-003-beth} is the \(L_{\omega_1\omega}\) form of
the first-order \cref{thm:p03-beth-definability}; its proof uses two
renamed copies of \(P\), as the proof of the first-order theorem does,
with the interpolation theorem in place of Robinson's separation
argument.
\Cref{thm:p02-w02c-p02-003-integers-second} applies
\cref{thm:godel-second} to an arithmetic theory obtained from \(T\)
through the translation of part~1 of
\cref{thm:p02-w02c-p02-003-integers}.  Its provability formula
(\cref{def:p02-w02c-p02-003-sharp-formula}) is the selected \(\PA\) formula
of \cref{def:p05-selected-direct-pa-graph} with one more axiom test, and
\cref{thm:p02-w02c-p02-003-sharp-provability} obtains D1--D3 for it from
\cref{thm:p05-direct-pa-internal-mp,thm:p05-original-sigma-evidence,thm:p05-literal-graph-evidence};
\cref{thm:p02-w02c-p02-003-arbitrary-representative} shows why some such
selection is needed.
Barwise passes from \(\mathcal N\) to \(\mathbb Z\) by appeal to the
standard definition of \(\mathbb Z\) from \(\mathcal N\)
\cite[ch.~A.1, \S2, p.~16]{Barwise1977FirstOrder};
\cref{thm:p02-w02c-p02-003-integers} goes through a definition of
\(\mathcal N\) inside \(\mathbb Z\), which is the direction in which a
reduction carries undecidability from \(\mathcal N\) to \(\mathbb Z\).

\emph{Sources.}  Barwise, ch.~A.1 of the \emph{Handbook of Mathematical
Logic}: \S2 pp.~8--9, 12, 15 (the sentences (6)--(8), the Archimedean
axiom, (10)), p.~16 (properties of theories), \S3 p.~17 (languages), \S5
pp.~43--45 (5.3--5.7), with \S3 pp.~17--18 and \S5 p.~44 (5.5)
\cite{Barwise1977FirstOrder}; Barwise, ``Axioms for abstract model
theory'', Part~I, \S1 p.~222, \S\S3--10 pp.~228--244 (definitions 4.1, 4.2,
6.2, the axioms of \S7, \S8, 9.1, 9.3, 9.7, 9.9, 9.10), Part~II, \S1
p.~249, \S2 pp.~254--255, \S3 pp.~257--259 (3.1, 3.4), \S4 pp.~260--262
(4.1), \S5 p.~264 (P5) \cite{Barwise1974Axioms}; Makkai, ch.~A.7, \S1
p.~236, \S2 pp.~238--239 and 243, \S4 pp.~249--251 (negation normal form,
4.1--4.3), \S5 p.~256, \S8 pp.~271--274 (8.1, 8.2, 8.6, and the remark on
consistency properties), with \S2 pp.~239--240, 242--243 (2.5, 2.8,
transitive closure, 2.16, 2.17) \cite{Makkai1977}; Rabin, ch.~C.3, \S2
pp.~608--609 (Theorem~4; algebraically closed fields)
\cite{Rabin1977Decidable}; Davis p.~234 \cite{Davis1973}.
